\documentclass[11pt,oneside]{article}

\usepackage[a4paper,margin=27mm,headheight=15pt]{geometry}
\usepackage[T1]{fontenc}
\usepackage[utf8]{inputenc}
\IfFileExists{libertinus.sty}{\usepackage{libertinus}}{\usepackage{lmodern}}
\usepackage{microtype}
\usepackage{amsmath,amssymb,amsthm,mathtools,mathrsfs}
\usepackage{bm}
\usepackage{booktabs,longtable,array,tabularx}
\usepackage{graphicx}
\usepackage{pgfplots}
\usepackage{xcolor}
\usepackage{enumitem}
\usepackage{fancyhdr}
\usepackage{titlesec}
\usepackage{listings}
\usepackage{xurl}
\usepackage{hyperref}
\usepackage[nameinlink,noabbrev]{cleveref}

\definecolor{linkblue}{RGB}{25,75,135}
\definecolor{shadegray}{RGB}{246,247,249}
\definecolor{rulegray}{RGB}{195,201,208}
\pgfplotsset{compat=1.18}
\hypersetup{
  colorlinks=true,
  linkcolor=linkblue,
  citecolor=linkblue,
  urlcolor=linkblue,
  pdftitle={Local, Boundary, and Reciprocal Expansions of q-Analogues},
  pdfauthor={Research report for the ProveIt project},
  pdfsubject={Finite and asymptotic expansions of q-Pochhammer symbols and Gaussian coefficients},
  pdfkeywords={q-Pochhammer symbol, Gaussian coefficient, q-binomial, Bernoulli polynomial, root of unity, q-gamma}
}

\pagestyle{fancy}
\fancyhf{}
\fancyhead[L]{\small q-analogue expansions}
\fancyhead[R]{\small Research-frontier report}
\fancyfoot[C]{\thepage}
\renewcommand{\headrulewidth}{0.4pt}

\titleformat{\section}{\Large\bfseries}{\thesection}{0.7em}{}
\titleformat{\subsection}{\large\bfseries}{\thesubsection}{0.7em}{}
\titleformat{\subsubsection}{\normalsize\bfseries}{\thesubsubsection}{0.7em}{}
\setcounter{secnumdepth}{3}
\setcounter{tocdepth}{2}
\setlist[itemize]{leftmargin=2em,itemsep=0.25em,topsep=0.4em}
\setlist[enumerate]{leftmargin=2.3em,itemsep=0.25em,topsep=0.4em}

% Long unbreakable inline mathematics occasionally leaves a line with no legal
% break point.  A small emergency stretch lets TeX loosen such a line rather
% than let it protrude into the margin; it applies only to lines that would
% otherwise be overfull.
\setlength{\emergencystretch}{2em}

\newtheorem{theorem}{Theorem}[section]
\newtheorem{proposition}[theorem]{Proposition}
\newtheorem{lemma}[theorem]{Lemma}
\newtheorem{corollary}[theorem]{Corollary}
\newtheorem{conjecture}[theorem]{Conjecture}
\theoremstyle{definition}
\newtheorem{definition}[theorem]{Definition}
\newtheorem{algorithm}[theorem]{Algorithm}
\newtheorem{example}[theorem]{Example}
\theoremstyle{remark}
\newtheorem{remark}[theorem]{Remark}
\newtheorem{warning}[theorem]{Warning}

\newcommand{\R}{\mathbb R}
\newcommand{\C}{\mathbb C}
\newcommand{\Q}{\mathbb Q}
\newcommand{\N}{\mathbb N}
\newcommand{\Z}{\mathbb Z}
\newcommand{\Up}{\operatorname{up}}
\newcommand{\sinc}{\operatorname{sinc}}
\newcommand{\wt}{\operatorname{wt}_2}
\newcommand{\e}{\mathrm e}
\newcommand{\ii}{\mathrm i}
\newcommand{\dd}{\,\mathrm d}
\newcommand{\Li}{\operatorname{Li}}
\newcommand{\Var}{\operatorname{Var}}
\newcommand{\E}{\mathbb E}
\newcommand{\Prob}{\mathbb P}
\newcommand{\bigO}{\mathcal O}
\newcommand{\smallo}{o}
\newcommand{\supp}{\operatorname{supp}}
\newcommand{\sgn}{\operatorname{sgn}}
\newcommand{\lcm}{\operatorname{lcm}}
\newcommand{\vtwo}{v_2}
\newcommand{\qbinom}[3]{\genfrac{[}{]}{0pt}{}{#1}{#2}_{#3}}
\newcommand{\repo}{\nolinkurl{Analysis/FabiusFunction}}
\newcommand{\lean}[1]{%
  \ifmmode\text{\nolinkurl{#1}}\else\nolinkurl{#1}\fi}
\newcommand{\leanevaltwo}[1]{%
  \nolinkurl{integral_eval}\texttt{\textsubscript{2}}\nolinkurl{#1}}

\newenvironment{proofidea}{\par\smallskip\noindent\textit{Proof architecture.}\ }{\hfill$\square$\par\smallskip}
\newenvironment{boxedremark}{%
  \par\smallskip\noindent\begin{center}\begin{minipage}{0.94\textwidth}%
  \color{black}\hrule\vspace{0.6em}\small
}{%
  \vspace{0.6em}\hrule\end{minipage}\end{center}\smallskip
}

\lstdefinestyle{pseudo}{
  basicstyle=\ttfamily\small,
  backgroundcolor=\color{shadegray},
  frame=single,
  rulecolor=\color{rulegray},
  columns=fullflexible,
  keepspaces=true,
  showstringspaces=false,
  xleftmargin=0.7em,
  xrightmargin=0.7em,
  aboveskip=0.8em,
  belowskip=0.8em
}

% Document-specific notation used below, appended after the shared block.
\newcommand{\ord}{\operatorname{ord}}
\newcommand{\Res}{\operatorname{Res}}
\newcommand{\cothh}{\operatorname{coth}}
\newcommand{\qbinomq}[2]{\qbinom{#1}{#2}{q}}
\newcommand{\qbinombase}[3]{\qbinom{#1}{#2}{#3}}
\newcommand{\qmultinom}[2]{\genfrac{[}{]}{0pt}{}{#1}{#2}_{q}}
\newcommand{\poch}[3]{(#1;#2)_{#3}}
\newcommand{\qpoch}[2]{(#1;q)_{#2}}
\newcommand{\Bell}{\mathcal Y}
\newcommand{\Thetaop}{\vartheta}
\newcommand{\eps}{\varepsilon}
\newcommand{\defeq}{\mathrel{:=}}
\newcommand{\abs}[1]{\lvert#1\rvert}
\newcommand{\floor}[1]{\left\lfloor#1\right\rfloor}
\newcommand{\ceil}[1]{\left\lceil#1\right\rceil}
\allowdisplaybreaks

% The arrival used a book-style two-level hierarchy.  These aliases map that
% hierarchy onto the canonical article counters without changing content.
\let\reporttopsection\section
\let\reportsubsection\subsection
\let\chapter\reporttopsection
\let\section\reportsubsection
\providecommand{\frontmatter}{}
\providecommand{\mainmatter}{}
\providecommand{\backmatter}{}

\title{\textbf{Local, Boundary, and Reciprocal Expansions\\of $q$-Pochhammer Symbols, Gaussian Coefficients,\\and Related $q$-Analogues}\\[0.8em]
\large Expansions at $q=1$, $q=0$, roots of unity, and $q=\infty$}
\author{Research report for the ProveIt Fabius-function project}
\date{Audited repository snapshot: 30 August 2026}

\begin{document}
\hypersetup{pageanchor=false}
\maketitle
\hypersetup{pageanchor=true}
\frontmatter

\chapter*{Abstract}
\addcontentsline{toc}{section}{Abstract}
This report develops a unified calculus for expansions of finite and infinite
$q$-products and the $q$-analogues built from them.  Four qualitatively different
situations are separated from the outset: ordinary Taylor expansions of finite
$q$-polynomials; boundary asymptotics of infinite products as $q$ approaches the
unit circle; Laurent expansions at the point $q=\infty$ for finite objects; and
generalized power or transseries expansions when a parameter coalesces with a
zero of the product.  This distinction prevents several common but consequential
confusions, such as treating $(q;q)_\infty$ as analytic at $q=1$, or speaking of
an unrenormalized infinite product at $\abs q>1$.

For finite shifted factorials we give all-order expansions both with a fixed
parameter $a$ and in the classical scaling $a=q^\alpha$.  For Gaussian and
multinomial coefficients, the logarithmic coefficients are expressed by
Bernoulli numbers and Faulhaber differences, while complete exponential Bell
polynomials recover every ordinary coefficient.  The centered Gaussian
coefficient is shown to be an even function of $t=-\log q$, and its coefficients
are identified with exact cumulants of the inversion statistic.  Around $q=0$
the coefficients are weighted restricted-partition numbers; reciprocity then
transfers the entire expansion to $q=\infty$ and distinguishes the positive and
negative real rays.  At $q=-1$ we obtain not only the familiar value of a
Gaussian coefficient but its complete first jet, including a closed derivative
formula in the vanishing parity case.  At a general root of unity, cyclotomic
factorization gives the zero order and a Bell-polynomial algorithm for all
Taylor coefficients.

For infinite products, the fixed-parameter $q\to1$ expansion is written in
polylogarithms, and the coalescing specialization $a=q^x$ is derived by Mellin
inversion.  This yields explicit all-order expansions of the $q$-gamma and
$q$-beta functions.  Dedekind-eta modularity supplies exponentially improved
formulae at $q=1$ and $q=-1$.  A radial expansion at every primitive root of
unity is derived directly from the Lambert series; its constant term is a
cyclic quantum dilogarithm.  Further sections treat the two standard
$q$-exponentials, the $q$-difference operator, basic hypergeometric series as
differential-operator deformations of ordinary hypergeometric series, and a
uniform double-scaling regime
\(
q=\exp(-\tau/N),\ k\sim\alpha N
\)
for Gaussian and multinomial coefficients.  The latter has an all-order
expansion in odd powers of $N^{-1}$ whose coefficients are elementary
polylogarithmic expressions.

The report was prepared as a self-contained companion to the ProveIt monograph
on $q$-Pochhammer symbols and $q$-binomial
coefficients~\cite{ReshetnikovMonograph}.  Several finite-algebra and
$q\to1$ statements below are rederivations of results already in that
monograph; the remaining report-level derivations are not thereby formalized
in Lean.  A companion Python program performs exact symbolic checks in finite
ranges and four high-precision spot checks.  Those checks support the printed
formulae but do not prove the general analytic remainder statements.

\chapter*{Principal results and status}
\addcontentsline{toc}{section}{Principal results and status}
\begin{boxedremark}
\textbf{Audit and provenance boundary.}
The extracted arrival was verified byte-for-byte against the five-entry
\path{MANIFEST.sha256} before normalization; that immutable file remains the
historical arrival ledger and has SHA-256
\texttt{17307793f182da4757f0344da4d7aec57b4f92fb8483005e92ea7b3c244a18cb}.
It is not a checksum ledger for the normalized files.  The separate
\path{SHA256SUMS.txt} records the current package.

In this report, ``proved here'' means that the displayed argument is a paper
proof in this semi-formalized frontier document.  It does \emph{not} mean that
an exact Lean declaration exists.  Only the declarations named in the
crosswalk below are claimed as live formal counterparts.  Likewise,
``companion-only'' means absent from the inspected monograph snapshot, not new
to the mathematical literature; no worldwide priority claim is made.
\end{boxedremark}

The following guide separates inherited baseline, report-level derivation,
asymptotic status, computation, and proposed research.

\begin{longtable}{@{}>{\raggedright\arraybackslash}p{0.22\textwidth}
                  >{\raggedright\arraybackslash}p{0.49\textwidth}
                  >{\raggedright\arraybackslash}p{0.20\textwidth}@{}}
\toprule
Topic & Main output & Status\\
\midrule
\endhead
Finite $q$-products at regular points & Bell-polynomial Taylor formula from logarithmic derivatives; resonant factor extraction & Proved here; not Lean as a jet theorem\\
$q\to1$, finite products & Fixed-$a$ coefficients and scaled $a=q^\alpha$ Bernoulli-polynomial coefficients & Mixed: companion derivation / monograph baseline\\
Gaussian coefficients & Logarithmic coefficients, centered expansion, ordinary Taylor coefficients, cumulants & Monograph baseline, extended here; not Lean\\
$q\to0$ & Weighted distinct-partition and rectangular-partition expansions & Monograph baseline / proved here\\
$q\to\infty$ & Reciprocal Laurent expansions and the two real rays & Baseline; base reversal is Lean-backed\\
$q=-1$ & Known value and the first derivative for every $n,k$ & Value is baseline; jet proved here, not Lean\\
General roots of unity & Cyclotomic order and carry criterion; Bell-polynomial Taylor jets & Baseline order/value; jets proved here, not Lean\\
Infinite product near $q=1$ & Fixed-$a$ polylogarithmic and coalescing Mellin expansions & Classical; asymptotic paper proofs; not Lean\\
$q$-gamma and $q$-beta & Explicit all-order logarithmic expansions & Asymptotic paper proofs; not Lean\\
Root-of-unity radial limit & Cyclic-dilogarithm factor and coefficient sums & Asymptotic paper proof plus one spot check\\
Double scaling & Expansion for $q=\e^{-\tau/N}$ and $k=\alpha N$ & Paper proof plus one spot check; not Lean\\
Research programme & Uniform resonant scaling, higher cyclotomic jets, resurgence, formal verification & Open / proposed\\
\bottomrule
\end{longtable}

\section*{Live-repository crosswalk}
\begin{tabularx}{\textwidth}{@{}>{\raggedright\arraybackslash}p{0.23\textwidth}X@{}}
\toprule
Layer & Audit finding\\
\midrule
Live Lean finite algebra &
\lean{FabiusFunction.FiniteQBinomialCore} proves the Pascal recursion,
symmetry, the denominator-free factorial relation, and the finite
$q$-binomial theorem.  \lean{FabiusFunction.GaussianBinomialAtOne} proves
evaluation at $q=1$.  \lean{FabiusFunction.QPochhammerElementaryIdentities}
proves finite base reversal.\\
Existing monograph &
The inspected monograph already contains the Gaussian Bernoulli expansion and
centering (\texttt{thm:gaussian-bernoulli}), rectangular-partition and
reciprocity results, cyclotomic factorization, the $q$-Lucas value at roots of
unity, and the value at $q=-1$.  Reappearances here are baseline or
rederivations, not novelty claims.\\
Unformalized report layer &
The general logarithmic/Bell jets, the closed derivative at $q=-1$, higher
cyclotomic jets, infinite-product and $q$-gamma asymptotics, radial
root-of-unity coefficients, hypergeometric operator deformation, and
double-scaling expansion have no exact proved Lean counterparts located in
the live source.  They remain research-frontier mathematics even when a paper
proof is supplied.\\
Computation &
The script checks the $q=1$ four-jet and $q=-1$ first jet for
$0\le k\le n\le20$, and performs four numerical evaluations at one parameter
point each.  It does not certify uniformity, optimal remainders, novelty, or
the two proposed research statements.\\
\bottomrule
\end{tabularx}

\tableofcontents

\mainmatter

\chapter{What kind of expansion is meant?}\label{chap:taxonomy}

\section{Four analytically different regimes}

The phrase ``expand at $q=q_0$'' hides several inequivalent operations.
Throughout the report they are kept separate.

\begin{enumerate}[label=\textup{(\roman*)}]
\item A finite shifted factorial $\qpoch an$ and a Gaussian coefficient
      $\qbinomq nk$ are polynomials in $q$.  They therefore possess ordinary,
      globally convergent Taylor series at every finite $q_0$.  The logarithm
      of such a polynomial is only local and is unavailable at a zero.
\item An infinite product $\qpoch a\infty$ is naturally defined for
      $\abs q<1$.  The points $q=1$ and the other points of the unit circle are
      boundary points, not ordinary Taylor centers.  The appropriate objects
      are sectorial asymptotic expansions, often with exponentials such as
      $\exp(-\Li_2(a)/t)$ and, in resonant cases, powers and logarithms of
      $t=-\log q$.
\item A finite polynomial has a Laurent expansion at the point $q=\infty$ on
      the Riemann sphere.  Reciprocity makes this expansion a reversed copy of
      its Taylor series at $q=0$.  The symbols $+\infty$ and $-\infty$ describe
      two real approaches to the same complex point and differ only through
      the phases of the inverse powers.
\item When a parameter depends on $q$ and approaches a zero of a product, the
      expansion generally leaves the category of ordinary powers.  Examples
      are $a=q^x$ as $q\to1$ and $\Gamma_q(x)$ as $q\to0$ for nonintegral $x$.
      Powers $q^x$, logarithms, and exponentially small corrections then form
      the natural asymptotic scale.
\end{enumerate}

This taxonomy is not merely terminological.  It determines which algebraic
operations are legitimate, what the remainder means, and whether an expansion
can be continued through the proposed center.

\section{Notation}

For $n\in\N_0$,
\[
 \qpoch an=\prod_{j=0}^{n-1}(1-aq^j),
 \qquad
 \qpoch a\infty=\prod_{j=0}^{\infty}(1-aq^j)\quad(\abs q<1),
\]
and
\[
 \qbinomq nk=\frac{(q;q)_n}{(q;q)_k(q;q)_{n-k}}
 \qquad(0\le k\le n),
\]
where the quotient is understood as its polynomial continuation.  We write
\[
 S_r(N)=\sum_{j=1}^{N}j^r,
 \qquad
 \Delta_r(n,k)=S_r(n)-S_r(k)-S_r(n-k),
 \qquad
 D=k(n-k).
\]
The Bernoulli numbers and polynomials use the convention
\[
 \frac{z\e^{xz}}{\e^z-1}
 =\sum_{r=0}^{\infty}B_r(x)\frac{z^r}{r!},
 \qquad B_r=B_r(0),
\]
so $B_1=-1/2$.  The polylogarithm is
\(
\Li_s(z)=\sum_{m\ge1}z^m/m^s
\)
in its initial disk and by analytic continuation elsewhere.

Whenever the logarithm of an infinite product occurs, it denotes the branch
defined by its absolutely convergent Lambert series and continued from
$a=0$ inside the stated simply connected parameter domain.  This convention
also fixes the logarithm of the cyclic product in the radial formulas; without
it, identities of logarithms would be determined only modulo $2\pi\ii$.

Two local coordinates will recur:
\[
 q=\e^{-t},\qquad t=-\log q,
 \qquad\text{and}\qquad h=q-1.
\]
The logarithmic coordinate $t$ exposes reciprocity and Bernoulli structure;
$h$ gives the literal Taylor expansion in $q$.

\section{Bell polynomials and conversion of coordinates}

\begin{definition}[Complete exponential Bell polynomials]
The polynomials $\Bell_r(x_1,\ldots,x_r)$ are defined by
\[
 \exp\left(\sum_{j\ge1}x_j\frac{z^j}{j!}\right)
 =\sum_{r\ge0}\Bell_r(x_1,\ldots,x_r)\frac{z^r}{r!}.
\]
Thus $\Bell_0=1$, $\Bell_1=x_1$,
$\Bell_2=x_1^2+x_2$, and
$\Bell_3=x_1^3+3x_1x_2+x_3$.
\end{definition}

\begin{lemma}[Exponentiation of a logarithmic expansion]\label{lem:bell-exp}
If
\[
 \log\frac{F(z)}{F(0)}
 =\sum_{j\ge1}\kappa_j\frac{z^j}{j!},
 \qquad F(0)\ne0,
\]
then
\[
 F(z)=F(0)\sum_{r\ge0}
 \Bell_r(\kappa_1,\ldots,\kappa_r)\frac{z^r}{r!}.
\]
The identity is formal; whenever the logarithmic series converges, the resulting
series converges in the same local disk.
\end{lemma}

\begin{proof}
This is the defining generating function of the complete exponential Bell
polynomials after substituting $x_j=\kappa_j$.
\end{proof}

Let $\Thetaop=q\frac{\dd}{\dd q}$.  Since $q=\e^{-t}$,
\[
 \frac{\dd}{\dd t}F(\e^{-t})=-\Thetaop F(q).
\]
Ordinary derivatives at $q=1$ can be recovered from logarithmic-coordinate
derivatives by
\begin{equation}\label{eq:stirling-coordinate}
 F^{(r)}(1)
 =\sum_{j=0}^{r}s(r,j)(-1)^j
 \left.\frac{\dd^j}{\dd t^j}F(\e^{-t})\right|_{t=0},
\end{equation}
where $s(r,j)$ are the signed Stirling numbers of the first kind.  Indeed,
$q^r\frac{\dd^r}{\dd q^r}=\Thetaop(\Thetaop-1)\cdots(\Thetaop-r+1)$.
Formula~\eqref{eq:stirling-coordinate}, together with
\cref{lem:bell-exp}, is an all-order converter between a compact cumulant-like
expansion in $t$ and a literal Taylor series in $q-1$.

\chapter{Finite products at an arbitrary center}\label{chap:regular-center}

\section{The universal regular-point formula}

\begin{theorem}[Taylor expansion from logarithmic derivatives]\label{thm:regular-bell}
Let $P(q)$ be holomorphic near $q_0$ and suppose $P(q_0)\ne0$.  Put
\[
 \lambda_r(q_0)=
 \left.\frac{\dd^r}{\dd q^r}\log P(q)\right|_{q=q_0}.
\]
Then
\[
 P(q_0+h)=P(q_0)
 \sum_{r\ge0}\frac{\Bell_r(\lambda_1,\ldots,\lambda_r)}{r!}h^r.
\]
For a polynomial the Taylor series terminates, although the logarithmic series
itself generally does not.
\end{theorem}

\begin{proof}
Apply \cref{lem:bell-exp} to $F(h)=P(q_0+h)$.
\end{proof}

For the finite shifted factorial, logarithmic differentiation may be performed
factor by factor:
\begin{equation}\label{eq:finite-poc-logder}
 \frac{\dd}{\dd q}\log\qpoch an
 =-\sum_{j=1}^{n-1}\frac{ajq^{j-1}}{1-aq^j}.
\end{equation}
Higher derivatives are rational functions with possible poles only at zeros of
the product.  A useful diagonal form follows from $\Thetaop=q\dd/\dd q$:
\begin{equation}\label{eq:euler-logder}
 \Thetaop^r\log\qpoch an
 =-\sum_{m\ge1}a^m m^{r-1}
   \sum_{j=0}^{n-1}j^r q^{jm}
 \qquad(r\ge1),
\end{equation}
initially when every $\abs{aq^j}<1$ and elsewhere by analytic continuation.
The finite geometric power sum in~\eqref{eq:euler-logder} can be reduced to a
rational function of $q^m$ or to Eulerian polynomials.

\section{Resonant zeros of a finite shifted factorial}

At a zero, the logarithmic expansion must be replaced by a factored expansion.
The zero multiplicity is completely elementary but becomes nontrivial at roots
of unity, where several factors may vanish simultaneously.

\begin{theorem}[Local leading term at a resonant point]\label{thm:poc-resonance}
Fix $n\ge1$, $q_0\ne0$, and $a\ne1$.  Let
\[
 J=\{j\in\{1,\ldots,n-1\}: aq_0^j=1\},
 \qquad r=\abs J.
\]
Then
\[
 \qpoch an=C(a,q_0,n)(q-q_0)^r\bigl(1+\bigO(q-q_0)\bigr),
\]
where
\begin{equation}\label{eq:poc-resonant-leading}
 C(a,q_0,n)=(-1)^r q_0^{-r}
 \left(\prod_{j\in J}j\right)
 \left(\prod_{\substack{0\le j<n\\j\notin J}}(1-aq_0^j)\right).
\end{equation}
After extracting the displayed power, all further coefficients follow from
\cref{thm:regular-bell} applied to the nonvanishing quotient.
\end{theorem}

\begin{proof}
For $j\in J$,
\[
 1-aq^j=-ajq_0^{j-1}(q-q_0)+\bigO((q-q_0)^2)
        =-\frac{j}{q_0}(q-q_0)+\bigO((q-q_0)^2).
\]
Multiplying these $r$ linear factors and evaluating every other factor at
$q_0$ gives~\eqref{eq:poc-resonant-leading}.  The quotient by
$(q-q_0)^r$ is holomorphic and nonzero at $q_0$, so the regular-point theorem
applies.
\end{proof}

\begin{remark}
If $a=1$, the factor with $j=0$ vanishes identically and $\qpoch 1n=0$ for all
$q$ and all $n\ge1$.  The important object $(q;q)_n$ is different: its first
parameter varies with $q$ and is treated by the scaled analysis in
\cref{chap:qone-finite}.
\end{remark}

\section{Gaussian coefficients through cyclotomic factors}

The polynomial factorization
\begin{equation}\label{eq:cyclotomic-factor}
 \qbinomq nk=\prod_{d=1}^{n}\Phi_d(q)^{e_d(n,k)},
 \qquad
 e_d(n,k)=\floor{\frac nd}-\floor{\frac kd}
          -\floor{\frac{n-k}{d}},
\end{equation}
has $e_d(n,k)\in\{0,1\}$.  It is therefore an optimal local representation at
roots of unity.

\begin{proposition}[Regular Taylor jets of a Gaussian polynomial]\label{prop:gaussian-regular-jet}
If $G_{n,k}(q)=\qbinomq nk$ and $G_{n,k}(q_0)\ne0$, then
\[
 \left(\log G_{n,k}\right)^{(r)}(q_0)
 =\sum_{d:e_d(n,k)=1}
   \left(\log\Phi_d\right)^{(r)}(q_0).
\]
Consequently every Taylor coefficient at $q_0$ is an explicit Bell polynomial
in logarithmic derivatives of the selected cyclotomic polynomials.
\end{proposition}

\begin{proof}
Take a local logarithm of~\eqref{eq:cyclotomic-factor} and apply
\cref{thm:regular-bell}.
\end{proof}

This construction is denominator-free and remains meaningful when the usual
factorial quotient is of the indeterminate form $0/0$.

\chapter{Finite products as \texorpdfstring{$q\to1$}{q tends to 1}}\label{chap:qone-finite}

\section{A fixed first parameter}

The source monograph emphasizes the scaled factors that build Gaussian
coefficients.  A complementary expansion keeps $a$ fixed while $q$ moves.  It
is useful for perturbing terminating basic hypergeometric identities whose
parameters are not integral powers of $q$.

\begin{theorem}[Fixed-$a$ finite shifted factorial]\label{thm:fixed-a-finite}
Let $n\ge1$, $q=\e^{-t}$, and initially assume $\abs a<1$.  Then, near $t=0$,
\begin{equation}\label{eq:fixed-a-finite-log}
 \log\qpoch a n
 =n\log(1-a)+
 \sum_{r\ge1}(-1)^{r+1}S_r(n-1)\Li_{1-r}(a)\frac{t^r}{r!}.
\end{equation}
The identity extends analytically in $a$ to every local domain avoiding the
zeros of $\qpoch an$ at $q=1$.  The ordinary coefficients of $\qpoch an$ are
the Bell polynomials in
\[
 \kappa_r=(-1)^{r+1}S_r(n-1)\Li_{1-r}(a).
\]
\end{theorem}

\begin{proof}
For $\abs a<1$,
\[
 \log(1-a\e^{-jt})
 =-\sum_{m\ge1}\frac{a^m}{m}\e^{-jmt}.
\]
Expanding the exponential and summing over $0\le j<n$ gives
\begin{align*}
 \log\qpoch an
 &=-\sum_{m\ge1}\frac{a^m}{m}
   \sum_{j=0}^{n-1}\sum_{r\ge0}\frac{(-jmt)^r}{r!}\\
 &=n\log(1-a)+
   \sum_{r\ge1}(-1)^{r+1}
   \left(\sum_{j=0}^{n-1}j^r\right)
   \left(\sum_{m\ge1}a^m m^{r-1}\right)\frac{t^r}{r!},
\end{align*}
which is~\eqref{eq:fixed-a-finite-log}.  The sum is finite in $j$, so the
result is a local analytic identity.  Analytic continuation in $a$ and
\cref{lem:bell-exp} finish the proof.
\end{proof}

The first terms are
\begin{align}
 \log\frac{\qpoch an}{(1-a)^n}
 &=\frac{n(n-1)}2\frac{a}{1-a}t
 -\frac{n(n-1)(2n-1)}{12}\frac{a}{(1-a)^2}t^2
 \notag\\
 &\quad+\frac{n^2(n-1)^2}{24}
 \frac{a(1+a)}{(1-a)^3}t^3+\bigO(t^4).
 \label{eq:fixed-a-first}
\end{align}
These terms already show why the fixed-$a$ and $a=q^\alpha$ limits must not be
interchanged: the rational functions of $a$ in~\eqref{eq:fixed-a-first} become
singular as $a\to1$.

\section{The classical scaling \texorpdfstring{$a=q^\alpha$}{a = q-alpha}}

\begin{theorem}[Finite scaled shifted factorial]\label{thm:scaled-finite}
Let $q=\e^{-t}$ and suppose none of
$\alpha,\alpha+1,\ldots,\alpha+n-1$ is zero.  Define
\[
 A_r(\alpha,n)=\sum_{j=0}^{n-1}(\alpha+j)^r
 =\frac{B_{r+1}(\alpha+n)-B_{r+1}(\alpha)}{r+1}.
\]
Then
\begin{equation}\label{eq:scaled-finite-log}
 \poch{q^\alpha}{q}{n}
 =t^n(\alpha)_n
 \exp\left(
 -\frac{A_1(\alpha,n)}2t
 +\sum_{r\ge1}\frac{B_{2r}A_{2r}(\alpha,n)}
 {2r(2r)!}t^{2r}
 \right).
\end{equation}
For fixed $\alpha,n$ the logarithmic series converges in every disk satisfying
$\abs{(\alpha+j)t}<2\pi$ for $0\le j<n$.  As a formal expansion it remains
valid after algebraic continuation, provided resonant zero factors are first
extracted.
\end{theorem}

\begin{proof}
The elementary identity
\begin{equation}\label{eq:basic-one-factor}
 1-\e^{-x}
 =x\exp\left(-\frac x2+
   \sum_{r\ge1}\frac{B_{2r}}{2r(2r)!}x^{2r}\right)
\end{equation}
holds for $\abs x<2\pi$.  Apply it to
$x=(\alpha+j)t$ and multiply over $0\le j<n$.  The product of the linear
factors is $t^n(\alpha)_n$, and the sums of powers are precisely
$A_r(\alpha,n)$.
\end{proof}

\begin{example}[The ordinary $q$-factorial]
Since
\(
[n]_q!=(q;q)_n/(1-q)^n
\), one obtains
\begin{equation}\label{eq:qfactorial-qone}
 \log\frac{[n]_{\e^{-t}}!}{n!}
 =-\frac{n(n-1)}4t+
 \sum_{r\ge1}\frac{B_{2r}}{2r(2r)!}
 \bigl(S_{2r}(n)-n\bigr)t^{2r}.
\end{equation}
For a single $q$-number,
\begin{equation}\label{eq:qnumber-qone}
 [x]_{\e^{-t}}
 =x\exp\left(
 -\frac{x-1}{2}t+
 \sum_{r\ge1}\frac{B_{2r}(x^{2r}-1)}{2r(2r)!}t^{2r}
 \right).
\end{equation}
The same formula is meaningful with $x$ replaced by an operator, a fact used in
\cref{chap:other-q}.
\end{example}

\chapter{Gaussian and multinomial coefficients near \texorpdfstring{$q=1$}{q = 1}}\label{chap:gaussian-qone}

\section{All logarithmic coefficients}

The basic Bernoulli formula in the source monograph is recalled here only to
make the later coefficient machinery self-contained.

\begin{theorem}[Complete Gaussian expansion]\label{thm:gaussian-complete}
Let $q=\e^{-t}$, $0\le k\le n$, and $D=k(n-k)$.  Then
\begin{equation}\label{eq:gaussian-log-all}
 \log\frac{\qbinombase nk{\e^{-t}}}{\binom nk}
 =-\frac D2t+
 \sum_{r\ge1}\frac{B_{2r}\Delta_{2r}(n,k)}{2r(2r)!}t^{2r}.
\end{equation}
For fixed $n,k$ this converges when $\abs t<2\pi/n$.  If
\[
 \kappa_1=-\frac D2,
 \qquad
 \kappa_{2r}=\frac{B_{2r}}{2r}\Delta_{2r}(n,k),
 \qquad
 \kappa_{2r+1}=0\quad(r\ge1),
\]
then the literal coefficient expansion is
\begin{equation}\label{eq:gaussian-bell}
 \qbinombase nk{\e^{-t}}
 =\binom nk\sum_{j\ge0}
 \Bell_j(\kappa_1,\ldots,\kappa_j)\frac{t^j}{j!}.
\end{equation}
\end{theorem}

\begin{proof}
Use the product
\[
 \qbinombase nk{\e^{-t}}
 =\prod_{j=1}^{k}
 \frac{1-\e^{-(n-k+j)t}}{1-\e^{-jt}},
\]
apply~\eqref{eq:basic-one-factor}, and collect power sums.  The linear
difference is $D$ and the even difference is
$\Delta_{2r}(n,k)$.  Bell exponentiation gives
\eqref{eq:gaussian-bell}.
\end{proof}

\section{Centering, explicit coefficients, and symmetry}

Reciprocity gives
\[
 \qbinombase nk{\e^{-t}}
 =\e^{-Dt}\qbinombase nk{\e^{t}}.
\]
Therefore
\begin{equation}\label{eq:centered-gaussian}
 C_{n,k}(t)=
 \e^{Dt/2}\frac{\qbinombase nk{\e^{-t}}}{\binom nk}
\end{equation}
is an even analytic function.  Put
\[
 c_{2r}=\frac{B_{2r}\Delta_{2r}}{2r(2r)!}.
\]
Then
\begin{align}
 C_{n,k}(t)
 &=1+c_2t^2+
   \left(c_4+\frac{c_2^2}{2}\right)t^4
 \notag\\
 &\quad+
   \left(c_6+c_2c_4+\frac{c_2^3}{6}\right)t^6
   +\bigO(t^8).                                      \label{eq:centered-through-six}
\end{align}
The needed Faulhaber differences factor as
\begin{align}
 \Delta_2&=D(n+1),\label{eq:delta2}\\
 \Delta_4&=D(n+1)\bigl(n(n+1)-D\bigr),\label{eq:delta4}\\
 \Delta_6&=\frac{D(n+1)}2
 \bigl(2D^2-4Dn^2-5Dn+2n^4+4n^3+n^2-n\bigr).
 \label{eq:delta6}
\end{align}
Consequently
\[
 c_2=\frac{D(n+1)}{24},
 \qquad
 c_4=-\frac{D(n+1)(n(n+1)-D)}{2880},
 \qquad
 c_6=\frac{\Delta_6}{181440}.
\]
All higher $\Delta_{2r}$ are symmetric polynomials in $k$ and $n-k$, hence
polynomials in $n$ and $D$.

\begin{corollary}[Literal Taylor series in $h=q-1$]\label{cor:qminusone-coordinate}
As $h\to0$,
\begin{align}
 \frac{\qbinombase nk{1+h}}{\binom nk}
 &=1+\frac D2h+
 \frac{D(3D+n-5)}{24}h^2
 \notag\\
 &\quad+
 \frac{D(D-2)(D+n-3)}{48}h^3+\bigO(h^4).
 \label{eq:gaussian-h-first}
\end{align}
Every later coefficient is obtained mechanically from
\eqref{eq:gaussian-log-all} by $t=-\log(1+h)$ and Bell expansion, or by the
Stirling conversion~\eqref{eq:stirling-coordinate}.
\end{corollary}

\begin{proof}
Substitute
$t=-h+h^2/2-h^3/3+\bigO(h^4)$ into
\eqref{eq:gaussian-log-all} and exponentiate.
\end{proof}

\section{Cumulants of inversions}

Let $X_{n,k}$ be the inversion number of a uniformly random binary word with
$k$ ones and $n-k$ zeros.  Its probability generating function is
\[
 \mathbb E(q^{X_{n,k}})=\frac{\qbinomq nk}{\binom nk}.
\]
Putting $q=\e^s$ in~\eqref{eq:gaussian-log-all} therefore gives all cumulants.

\begin{corollary}[Exact inversion cumulants]\label{cor:inversion-cumulants}
The odd centered cumulants vanish, and
\begin{align*}
 \operatorname{cum}_1(X_{n,k})&=\frac D2,\\
 \operatorname{cum}_{2r}(X_{n,k})
 &=\frac{B_{2r}}{2r}\Delta_{2r}(n,k)\qquad(r\ge1),\\
 \operatorname{cum}_{2r+1}(X_{n,k})&=0\qquad(r\ge1).
\end{align*}
In particular,
\[
 \operatorname{Var}(X_{n,k})=\frac{D(n+1)}{12},
\]
and the fourth cumulant is
\[
 -\frac{D(n+1)(n(n+1)-D)}{120}.
\]
\end{corollary}

The evenness in~\eqref{eq:centered-gaussian} is thus exactly the symmetry of
the inversion distribution about $D/2$, while Bernoulli numbers encode its
entire cumulant hierarchy.

\section{Multinomial extension}

Let $n_1+\cdots+n_s=n$ and write
\[
 \qmultinom{n}{n_1,\ldots,n_s}
 =\frac{(q;q)_n}{(q;q)_{n_1}\cdots(q;q)_{n_s}}.
\]
Define
\[
 \Delta_r(\bm n)=S_r(n)-\sum_{i=1}^{s}S_r(n_i),
 \qquad
 D(\bm n)=\Delta_1(\bm n)=\sum_{i<j}n_i n_j.
\]
The same proof gives
\begin{equation}\label{eq:qmult-qone}
 \log\frac{\qmultinom{n}{n_1,\ldots,n_s}\big|_{q=\e^{-t}}}
 {\binom{n}{n_1,\ldots,n_s}}
 =-\frac{D(\bm n)}2t+
 \sum_{r\ge1}\frac{B_{2r}\Delta_{2r}(\bm n)}{2r(2r)!}t^{2r}.
\end{equation}
Thus the centered $q$-multinomial is even, and
$B_{2r}\Delta_{2r}(\bm n)/(2r)$ is the $2r$th cumulant of the inversion
number of a uniformly random multiset permutation.

\chapter{Expansions at \texorpdfstring{$q=0$}{q = 0}}\label{chap:qzero}

\section{Finite and infinite shifted factorials}

For $n\ge1$,
\[
 \qpoch an=(1-a)\prod_{j=1}^{n-1}(1-aq^j).
\]
Every coefficient therefore has a direct restricted-partition meaning.

\begin{theorem}[Weighted distinct-partition expansion]\label{thm:qzero-poc}
For $r\ge0$ define
\[
 d_{n,r}(a)=
 \sum_{\substack{J\subseteq\{1,\ldots,n-1\}\\\sum_{j\in J}j=r}}
 (-a)^{\abs J}.
\]
Then
\begin{equation}\label{eq:qzero-poc-exact}
 \qpoch an=(1-a)\sum_{r=0}^{n(n-1)/2}d_{n,r}(a)q^r.
\end{equation}
For fixed $r$, $d_{n,r}(a)$ stabilizes once $n>r$ to
\[
 d_r(a)=\sum_{\substack{\lambda\vdash r\\\lambda\text{ has distinct parts}}}
 (-a)^{\ell(\lambda)},
\]
and hence
\begin{equation}\label{eq:qzero-poc-infinite}
 \qpoch a\infty=(1-a)\sum_{r\ge0}d_r(a)q^r.
\end{equation}
\end{theorem}

\begin{proof}
Expand the finite product by choosing either $1$ or $-aq^j$ from each
factor.  A choice of indices is a subset $J$ and contributes
$(-a)^{\abs J}q^{\sum J}$.  Factors with index at least $n$ cannot affect a
coefficient of degree below $n$, which proves stabilization and the infinite
identity.
\end{proof}

The beginning of the stable expansion is
\begin{align}
 \frac{\qpoch a\infty}{1-a}
 &=1-aq-aq^2+a(a-1)q^3+a(a-1)q^4
 \notag\\
 &\quad+a(2a-1)q^5-a(a-1)^2q^6
 -a(a^2-3a+1)q^7+\bigO(q^8).
 \label{eq:qzero-poc-first}
\end{align}
The finite product has the same terms through $q^R$ whenever $n>R$.

\begin{corollary}[Reciprocal product and unrestricted partitions]
For $a\ne1$,
\begin{equation}\label{eq:qzero-reciprocal-poc}
 \frac1{\qpoch a\infty}
 =\frac1{1-a}\sum_{r\ge0}
 \left(\sum_{\lambda\vdash r}a^{\ell(\lambda)}\right)q^r.
\end{equation}
In particular,
\begin{align*}
 \frac1{\qpoch a\infty}
 =\frac1{1-a}\bigl(&1+aq+a(a+1)q^2
 +a(a^2+a+1)q^3\\
 &+a(a^3+a^2+2a+1)q^4+\bigO(q^5)\bigr).
\end{align*}
\end{corollary}

\begin{proof}
Expand each $(1-aq^j)^{-1}$ geometrically.  The multiplicity chosen from the
$j$th factor is the multiplicity of the part $j$; the total power of $a$ is the
number of parts.
\end{proof}

Two familiar resonant specializations deserve separate mention:
\begin{align}
 (q;q)_\infty
 &=\sum_{m\in\Z}(-1)^m q^{m(3m-1)/2}
 =1-q-q^2+q^5+q^7-q^{12}-q^{15}+\cdots,\label{eq:pentagonal}\\
 \frac1{(q;q)_\infty}&=\sum_{r\ge0}p(r)q^r.\label{eq:partition-generating}
\end{align}
They do not follow by naively putting $a=1$ in
\eqref{eq:qzero-poc-infinite}, because in $(q;q)_\infty$ the first parameter
moves with $q$.

\section{Gaussian polynomials and rectangular partitions}

\begin{theorem}[Complete $q=0$ expansion of a Gaussian coefficient]
\label{thm:qzero-gaussian}
Let $D=k(n-k)$.  Then
\begin{equation}\label{eq:qzero-gaussian}
 \qbinomq nk=\sum_{r=0}^{D}p_{k,n-k}(r)q^r,
\end{equation}
where $p_{u,v}(r)$ is the number of partitions of $r$ fitting in a
$u\times v$ rectangle.  Consequently,
\[
 [q^r]\qbinomq nk=p(r)
 \qquad\text{whenever}\qquad r\le\min(k,n-k).
\]
The coefficients are palindromic:
$p_{k,n-k}(r)=p_{k,n-k}(D-r)$.
\end{theorem}

\begin{proof}
The standard Ferrers-diagram interpretation gives
\eqref{eq:qzero-gaussian}.  Every partition of $r$ has at most $r$ parts and
largest part at most $r$, proving the stable range.  Complementation inside the
rectangle proves palindromicity.
\end{proof}

Thus, if the rectangle is sufficiently large,
\[
 \qbinomq nk=1+q+2q^2+3q^3+5q^4+7q^5+11q^6+\cdots
\]
for as many initial terms as fit.  This observation is often more useful than
a quotient expansion: it gives coefficients exactly and certifies their
nonnegativity.

For a $q$-multinomial, the coefficient of $q^r$ counts multiset words of the
prescribed content and inversion number $r$.  Hence its constant term is one,
its degree is $D(\bm n)=\sum_{i<j}n_i n_j$, and reciprocity makes its
coefficient sequence palindromic.

\section{The small-\texorpdfstring{$q$}{q} expansion of \texorpdfstring{$\Gamma_q$}{q-gamma}}

For nonintegral $x$, $q^x$ prevents an ordinary Taylor series at zero.  The
right replacement is a convergent generalized-power expansion.

\begin{proposition}[Exact logarithmic expansion at $q=0$]\label{prop:qgamma-zero}
For $0<q<1$ and $x>0$,
\begin{align}
 \log\Gamma_q(x)
 =&\sum_{m\ge1}\frac{q^{mx}}{m(1-q^m)}
 -\sum_{m\ge1}\frac{q^m}{m(1-q^m)}
 -(1-x)\sum_{m\ge1}\frac{q^m}{m}.                 \label{eq:qgamma-zero-log}
\end{align}
All three series converge absolutely.  Expanding
$(1-q^m)^{-1}=\sum_{r\ge0}q^{mr}$ gives a convergent generalized series whose
logarithmic support lies in
\[
 \{m(x+r):m\ge1,r\ge0\}\cup\N.
\]
For rational $x$ this is a Puiseux series; for irrational $x$ it is naturally
a Hahn-type generalized power series.  In particular, the contributions from the two primitive generators of the
support begin as
\[
 \log\Gamma_q(x)=q^x-(2-x)q+\cdots.
\]
The omitted exponents lie in the displayed additive support and need not occur
in increasing order relative to both $x$ and $1$; collisions cause additional
cancellations at special integral values such as $x=1,2$.
\end{proposition}

\begin{proof}
Insert the product definition
\(
\Gamma_q(x)=(1-q)^{1-x}(q;q)_\infty/(q^x;q)_\infty
\)
and use
\(
\log(a;q)_\infty=-\sum_{m\ge1}a^m/(m(1-q^m)).
\)
Absolute convergence follows from $0<q<1$ and $x>0$.
\end{proof}

\chapter{The reciprocal point \texorpdfstring{$q=\infty$}{q = infinity} and the two real rays}\label{chap:infinity}

\section{Finite shifted factorials}

\begin{theorem}[Exact base reversal]\label{thm:poc-infinity}
For $a\ne0$ and $n\ge0$,
\begin{equation}\label{eq:poc-reciprocity}
 \qpoch an=(-a)^nq^{\binom n2}\poch{a^{-1}}{q^{-1}}{n}.
\end{equation}
Consequently, with $p=q^{-1}$,
\begin{align}
 \qpoch an
 &=(1-a)(-a)^{n-1}q^{\binom n2}
 \Bigl(1-a^{-1}p-a^{-1}p^2
 \notag\\
 &\hspace{7em}+(a^{-2}-a^{-1})p^3+\bigO(p^4)\Bigr),
 \label{eq:poc-infinity-first}
\end{align}
provided $n>3$; for smaller $n$ the series truncates in the evident way.
Every Laurent coefficient at infinity is obtained from the corresponding
$q=0$ coefficient by the substitution $a\mapsto a^{-1}$.
\end{theorem}

\begin{proof}
Factor $-aq^j$ from each factor $1-aq^j$ and multiply.  The second statement
is~\eqref{eq:qzero-poc-first} applied to the reversed product.
\end{proof}

For $q=R\to+\infty$, all powers $R^{-j}$ retain their signs.  For
$q=-R\to-\infty$, the coefficient of $R^{-j}$ acquires an additional factor
$(-1)^j$, and the leading phase is $(-1)^{\binom n2}$.  Thus $+\infty$ and
$-\infty$ are not distinct complex singularities, but they are distinct real
asymptotic rays.

The ordinary infinite product does not converge for $\abs q>1$.  There is,
however, a canonical renormalized finite-product limit.

\begin{corollary}[Renormalized product outside the unit disk]
If $\abs q>1$ and $a\ne0$, then
\begin{equation}\label{eq:outside-renormalization}
 \lim_{n\to\infty}
 \frac{\qpoch an}{(-a)^nq^{\binom n2}}
 =\poch{a^{-1}}{q^{-1}}{\infty}.
\end{equation}
\end{corollary}

\begin{proof}
Use~\eqref{eq:poc-reciprocity}; the product on the right converges because
$\abs{q^{-1}}<1$.
\end{proof}

\section{Gaussian and multinomial reciprocity}

\begin{theorem}[Complete Laurent expansion of a Gaussian coefficient]
\label{thm:gaussian-infinity}
With $D=k(n-k)$,
\begin{equation}\label{eq:gaussian-reciprocity}
 \qbinomq nk=q^D\qbinombase nk{q^{-1}}
 =q^D\sum_{r=0}^{D}p_{k,n-k}(r)q^{-r}.
\end{equation}
Thus the expansion at infinity begins
\[
 q^D\bigl(1+q^{-1}+2q^{-2}+3q^{-3}+\cdots\bigr)
\]
for as long as $r\le\min(k,n-k)$.  Along the negative real ray,
\begin{equation}\label{eq:gaussian-minus-infinity}
 \qbinombase nk{-R}
 =(-1)^D R^D\sum_{r=0}^{D}(-1)^r p_{k,n-k}(r)R^{-r}.
\end{equation}
\end{theorem}

\begin{proof}
The first identity follows either from the factorial quotient or from
palindromicity of the Gaussian polynomial.  Insert
\eqref{eq:qzero-gaussian}; then substitute $q=-R$.
\end{proof}

The identical argument gives, for a composition $\bm n$,
\[
 \qmultinom{n}{n_1,\ldots,n_s}
 =q^{D(\bm n)}
 \qmultinom{n}{n_1,\ldots,n_s}\bigg|_{q\mapsto q^{-1}}.
\]
The Laurent coefficients are the inversion numbers of multiset words read in
reverse order.

For $q$-numbers and factorials one has the exact relations
\[
 [x]_q=q^{x-1}[x]_{q^{-1}}
 \quad\text{for integral }x,
 \qquad
 [n]_q!=q^{\binom n2}[n]_{q^{-1}}!.
\]
These are the simplest instances of the same reciprocal principle.

\chapter{The point \texorpdfstring{$q=-1$}{q = -1}}\label{chap:qminusone}

\section{Finite shifted factorials}

Pairing even and odd exponents gives the exact value
\begin{equation}\label{eq:poc-minus-one-value}
 \qpoch a n\big|_{q=-1}
 =(1-a)^{\ceil{n/2}}(1+a)^{\floor{n/2}}.
\end{equation}
In particular, for every $m\geq0$,
\begin{equation}\label{eq:poc-minus-one-paired}
 (a;-1)_{2m}=(1-a^2)^m,
 \qquad
 (a;-1)_{2m+1}=(1-a^2)^m(1-a).
\end{equation}
Indeed, the factor indexed by an even $j$ is $1-a$, while the factor indexed
by an odd $j$ is $1+a$.  A product of length $2m$ contains $m$ factors of
each kind, so pairing them gives $((1-a)(1+a))^m=(1-a^2)^m$; a product of
length $2m+1$ has precisely one additional even-index factor.  This argument
is valid in an arbitrary commutative ring and is the proof formalized by
\path|finiteQPochhammerIn_neg_one_even| and
\path|finiteQPochhammerIn_neg_one_odd|.
When $a\ne\pm1$, the first derivative follows from
\eqref{eq:finite-poc-logder}.  Writing $n=2m$ or $2m+1$ gives
\begin{align}
 \left.\frac{\dd}{\dd q}\log\qpoch a{2m}\right|_{q=-1}
 &=\frac{am((2m-1)a-1)}{1-a^2},\label{eq:poc-minus-one-even-der}\\
 \left.\frac{\dd}{\dd q}\log\qpoch a{2m+1}\right|_{q=-1}
 &=\frac{am(1+(2m+1)a)}{1-a^2}.\label{eq:poc-minus-one-odd-der}
\end{align}
Higher derivatives are supplied by \cref{thm:regular-bell}.  If $a=-1$,
the odd-index factors vanish and
\begin{equation}\label{eq:a-minus-one-resonance}
 (-1;q)_n
 =2^{\ceil{n/2}}(2\floor{n/2}-1)!!
 (q+1)^{\floor{n/2}}\bigl(1+\bigO(q+1)\bigr),
\end{equation}
where $(-1)!!=1$.  The specialization $a=1$ is identically zero for $n\ge1$.

\section{\texorpdfstring{$q$}{q}-factorials}

The $q$-integer satisfies
\[
 [2r]_{-1}=0,
 \qquad [2r+1]_{-1}=1,
\]
and its first nonzero term is
\[
 [2r]_q=r(q+1)+\bigO((q+1)^2).
\]
It follows that
\begin{equation}\label{eq:qfactorial-minus-one}
 [n]_q!
 =\floor{n/2}!\,(q+1)^{\floor{n/2}}
 \bigl(1+\bigO(q+1)\bigr).
\end{equation}
This formula is a convenient way to evaluate factorial quotients at $-1$
without ever forming a numerical $0/0$.

\section{The complete first jet of a Gaussian coefficient}

\begin{theorem}[Value and derivative at $q=-1$]\label{thm:gaussian-minus-one-jet}
Assume $0\leq k\leq n$ and write
\[
 n=2a+\eps,
 \qquad k=2b+\delta,
 \qquad \eps,\delta\in\{0,1\}.
\]
If $(\eps,\delta)\ne(0,1)$, then
\begin{align}
 \qbinomq nk\bigg|_{q=-1}&=\binom ab,\label{eq:minus-one-value-nonzero}\\
 \left.\frac{\dd}{\dd q}\qbinomq nk\right|_{q=-1}
 &=-\frac{k(n-k)}2\binom ab.\label{eq:minus-one-der-nonzero}
\end{align}
If $n=2a$ and $k=2b+1$, then
\begin{align}
 \qbinomq{2a}{2b+1}\bigg|_{q=-1}&=0,\label{eq:minus-one-value-zero}\\
 \left.\frac{\dd}{\dd q}\qbinomq{2a}{2b+1}\right|_{q=-1}
 &=(a-b)\binom ab
 =a\binom{a-1}{b}.                                  \label{eq:minus-one-der-zero}
\end{align}
Thus every zero at $-1$ among these admissible Gaussian coefficients is
simple over $\Z$, and therefore after base change to any field of
characteristic zero.  After reduction modulo a prime, the displayed nonzero
integer slope can vanish, so simplicity in positive characteristic is a
separate question.
\end{theorem}

\begin{proof}
Here is a denominator-free proof of the values.  Work first over $\Z$, put
$G_{n,k}=\qbinombase{n}{k}{-1}$ and extend the Gaussian triangle by zero outside
$0\leq k\leq n$.  Reciprocity has the polynomial form
$\qbinomq{n}{k}=q^{k(n-k)}\qbinombase{n}{k}{q^{-1}}$.  In an even row and an
admissible odd column the exponent $k(n-k)$ is odd, so evaluation at $q=-1$ gives
$G_{2a,2b+1}=-G_{2a,2b+1}$; the value is an integer, hence it is zero.
Transport along the canonical homomorphism $\Z\to R$ gives the same vanishing
identity in every commutative ring $R$, including rings with $2$-torsion.

Now use the denominator-free alternate $q$-Pascal recurrence
\[
 \qbinomq{n+1}{k+1}
 =q^{k+1}\qbinomq{n}{k+1}+\qbinomq{n}{k}.
\]
Suppose the even row $2a$ has
$G_{2a,2b}=\binom ab$ and $G_{2a,2b+1}=0$.  The two intervening odd-row
entries are then
\[
 G_{2a+1,2b+1}=\binom ab,
 \qquad
 G_{2a+1,2b+2}=\binom a{b+1}.
\]
Applying the same recurrence once more gives
\[
 G_{2a+2,2b+2}
 =\binom a{b+1}+\binom ab
 =\binom{a+1}{b+1}.
\]
The boundary entries are $1$, so induction from row zero proves all four
parity cases, including the zero-extended cases above the row.  This is also
the $m=2$ case of $q$-Lucas and agrees with cancelling equal powers of $q+1$
in~\eqref{eq:qfactorial-minus-one}.

Assume first that the value is nonzero.  The degree is
$D=k(n-k)$, which is even in precisely these parity cases.  Reciprocity gives
$G(q)=q^DG(q^{-1})$ for $G(q)=\qbinomq nk$.  Differentiating and setting
$q=-1$ yields
\[
 G'(-1)=-DG(-1)-G'(-1),
\]
so $G'(-1)=-DG(-1)/2$, proving
\eqref{eq:minus-one-der-nonzero}.

Now let $n=2a$ and $k=2b+1$.  In
\[
 G(q)=\frac{(q;q)_{2a}}
 {(q;q)_{2b+1}(q;q)_{2a-2b-1}},
\]
the numerator has $a$ even factors, while the denominator has
$b+(a-b-1)=a-1$ even factors.  Hence one factor $q+1$ remains.  Moreover,
\[
 \frac{(q;q)_{2r}}{(q+1)^r}\longrightarrow 4^r r!,
 \qquad
 \frac{(q;q)_{2r+1}}{(q+1)^r}\longrightarrow 2^{2r+1}r!.
\]
Taking the quotient gives
\[
 \lim_{q\to-1}\frac{G(q)}{q+1}
 =\frac{a!}{b!(a-b-1)!}
 =(a-b)\binom ab,
\]
which is the derivative at the simple zero.
\end{proof}

\begin{remark}[Lean status of the second-root calculation]
The complete four-case value table is machine-checked over every commutative
ring.  Its four exact declarations are:
\begin{itemize}[nosep,leftmargin=2em]
\item \path|Fabius.gaussianBinomial_neg_one_even_even|;
\item \path|Fabius.gaussianBinomial_neg_one_even_odd_eq_zero|;
\item \path|Fabius.gaussianBinomial_neg_one_odd_even|;
\item \path|Fabius.gaussianBinomial_neg_one_odd_odd|.
\end{itemize}
The odd-column zero lives in \path|FabiusFunction.QBinomialReciprocity|, where
it follows from reciprocal symmetry.  The other three values and the two
product identities \path|Fabius.finiteQPochhammerIn_neg_one_even| and
\path|Fabius.finiteQPochhammerIn_neg_one_odd| live in
\path|FabiusFunction.GaussianBinomialAtNegOne|.  That module reuses the zero
theorem, then needs only $q$-Pascal, parity, ordinary Pascal, and factor pairing.
The first-jet formulas~\eqref{eq:minus-one-der-nonzero} and
\eqref{eq:minus-one-der-zero} remain the next formalization layer; this report
does not treat numerical checks as a substitute for that proof.  Their eventual
formalization requires a polynomial Gaussian-coefficient interface; simplicity
must be stated over the integers or in characteristic zero, because a nonzero
integral slope may vanish after reduction modulo a prime.
\end{remark}

\begin{remark}[Why the derivative formula is structurally natural]
In the nonvanishing cases, palindromicity alone fixes the derivative at $-1$.
In the vanishing case, parity forces odd degree and reciprocity only predicts
that the value must vanish; the uncancelled even factorial factor supplies the
missing slope.  The two mechanisms together determine the entire first jet.
\end{remark}

\chapter{General roots of unity}\label{chap:roots-finite}

\section{Cyclotomic order and the carry criterion}

Let $\zeta$ be a primitive $m$th root of unity.  From
\eqref{eq:cyclotomic-factor},
\begin{equation}\label{eq:root-order}
 \ord_{q=\zeta}\qbinomq nk=e_m(n,k)\in\{0,1\}.
\end{equation}
Write $n=am+r$ and $k=bm+s$ with $0\le r,s<m$.  A one-line floor calculation
shows
\begin{equation}\label{eq:carry-criterion}
 e_m(n,k)=1\quad\Longleftrightarrow\quad s>r.
\end{equation}
Thus the zero records exactly the carry that occurs when the residues of $k$
and $n-k$ are added in base $m$.

The $q$-Lucas value is
\begin{equation}\label{eq:q-lucas-root}
 \qbinomq nk\bigg|_{q=\zeta}
 =\binom ab\qbinombase rs\zeta,
\end{equation}
where the right side is zero when $s>r$.

\section{All local Taylor coefficients}

\begin{theorem}[Cyclotomic jet algorithm]\label{thm:cyclotomic-jets}
Let
\[
 G(q)=\qbinomq nk,
 \qquad e=e_m(n,k),
 \qquad H(q)=\frac{G(q)}{(q-\zeta)^e}.
\]
Then $H$ is holomorphic and nonzero at $\zeta$.  More explicitly,
\begin{equation}\label{eq:root-H}
 H(q)=
 \left(\frac{\Phi_m(q)}{q-\zeta}\right)^e
 \prod_{\substack{d\ne m\\e_d(n,k)=1}}\Phi_d(q).
\end{equation}
If
\(
\lambda_r=(\dd^r/\dd q^r)\log H(q)|_{q=\zeta}
\), then
\begin{equation}\label{eq:root-jet-bell}
 G(q)=(q-\zeta)^eH(\zeta)
 \sum_{r\ge0}\Bell_r(\lambda_1,\ldots,\lambda_r)
 \frac{(q-\zeta)^r}{r!}.
\end{equation}
In particular, if $e=1$,
\begin{equation}\label{eq:root-leading-derivative}
 G'(\zeta)=\Phi_m'(\zeta)
 \prod_{\substack{d\ne m\\e_d(n,k)=1}}\Phi_d(\zeta).
\end{equation}
\end{theorem}

\begin{proof}
Only $\Phi_m$ vanishes at a primitive $m$th root, and it has a simple zero
there.  Formula~\eqref{eq:root-H} is therefore holomorphic and nonzero.
Apply \cref{thm:regular-bell} to $H$ and restore the extracted factor.
\end{proof}

This is an exact finite algorithm: determine the selected cyclotomic factors
from floor functions, differentiate their logarithms symbolically or in the
number field $\Q(\zeta)$, and evaluate Bell polynomials.  It avoids unstable
cancellation in a quotient of factorials and works to arbitrary order.

\chapter{Infinite shifted factorials as \texorpdfstring{$q\to1$}{q tends to 1}}\label{chap:infinite-qone}

\section{Fixed \texorpdfstring{$a$}{a}: the polylogarithmic expansion}

For $\abs q<1$ and $\abs a<1$, the Lambert representation is
\begin{equation}\label{eq:lambert-log-poc}
 \log\qpoch a\infty
 =-\sum_{m\ge1}\frac{a^m}{m(1-q^m)}.
\end{equation}
It is the natural starting point because the singularity as $q\to1$ is
isolated in the elementary kernel $(1-q^m)^{-1}$.

\begin{theorem}[Fixed-parameter all-order expansion]\label{thm:infinite-fixed-a}
Let $q=\e^{-t}$ with $t\to0$ in a closed subsector of $\Re t>0$.  Uniformly for
$\abs a\le\rho<1$,
\begin{equation}\label{eq:infinite-fixed-a}
 \log\poch a{\e^{-t}}\infty
 \sim-\frac{\Li_2(a)}{t}+\frac12\log(1-a)
 -\sum_{r\ge1}\frac{B_{2r}}{(2r)!}
   \Li_{2-2r}(a)t^{2r-1}.
\end{equation}
Equivalently,
\begin{equation}\label{eq:infinite-fixed-a-product}
 \poch a{\e^{-t}}\infty
 \sim (1-a)^{1/2}\exp\left(-\frac{\Li_2(a)}t
 \right)
 \exp\left(-\sum_{r\ge1}\frac{B_{2r}}{(2r)!}
 \Li_{2-2r}(a)t^{2r-1}\right).
\end{equation}
After the dominant exponential and square-root factor are removed, every
ordinary coefficient follows from complete Bell polynomials.
\end{theorem}

\begin{proof}
Use
\[
 \frac1{1-\e^{-x}}
 =\frac1x+\frac12+
 \sum_{r=1}^{N-1}\frac{B_{2r}}{(2r)!}x^{2r-1}+R_N(x)
\]
in~\eqref{eq:lambert-log-poc}.  The first three kinds of terms sum to
$-\Li_2(a)/t$, $\tfrac12\log(1-a)$, and the displayed polylogarithms.
For the remainder, split the $m$-sum at $m\le\abs t^{-1}$.  On the first part,
Taylor's theorem and $\abs a\le\rho$ give
\(
\sum \rho^m m^{2N-2}\abs t^{2N-1}=\bigO(\abs t^{2N-1})
\).
The tail is exponentially small in $1/\abs t$ because of $\rho^m$.  This proves
the uniform Poincare expansion.
\end{proof}

The first correction is
\[
 -\frac{t}{12}\Li_0(a)
 =-\frac{t}{12}\frac{a}{1-a},
\]
and the next nonzero logarithmic correction is
\[
 +\frac{t^3}{720}\Li_{-2}(a)
 =\frac{t^3}{720}\frac{a(1+a)}{(1-a)^3}.
\]
The absence of even positive powers is a property of the fixed-$a$ regime; it
is destroyed when $a$ itself approaches one.

Formula~\eqref{eq:infinite-fixed-a} is classical in spirit and is closely
related to the complete expansions developed by McIntosh
\cite{McIntoshShifted}.  A recent gamma-product identity of Arabi Ardehali and
Rosengren gives an exact completion in which the asymptotic remainder is an
absolutely convergent product of dressed gamma factors
\cite{ArdehaliRosengren}.

\section{The coalescing parameter \texorpdfstring{$a=q^x$}{a = q to x}}

The substitution $a=\e^{-xt}$ lies outside every compact set
$\abs a\le\rho<1$.  The singularity of $\Li_2(a)$ at $a=1$ interacts with
$t\to0$, producing powers of $t$ and $\log t$ governed by the gamma function.

\begin{theorem}[Mellin expansion of $(q^x;q)_\infty$]
\label{thm:qx-infinite}
Let $\Re x>0$ and $q=\e^{-t}$, where $t\to0$ in a closed subsector of
$\abs{\arg t}<\pi/2$.  Then
\begin{align}
 \log\poch{q^x}{q}\infty
 \sim{}&-\frac{\pi^2}{6t}
 +\left(\frac12-x\right)\log t
 +\frac12\log(2\pi)-\log\Gamma(x)
 \notag\\
 &+\frac{B_2(x)}4t
 -\sum_{r\ge1}
 \frac{B_{2r}B_{2r+1}(x)}{2r(2r+1)!}t^{2r}.
 \label{eq:qx-infinite}
\end{align}
The expansion is locally uniform for $x$ in compact subsets of
$\Re x>0$.
\end{theorem}

\begin{proof}
For $\Re s>1$, Mellin inversion of the exponential in the Lambert series gives
\begin{equation}\label{eq:qx-mellin}
 \log\poch{\e^{-xt}}{\e^{-t}}\infty
 =-\frac1{2\pi\ii}\int_{c-\ii\infty}^{c+\ii\infty}
 \Gamma(s)\zeta(s+1)\zeta(s,x)t^{-s}\dd s,
 \qquad c>1.
\end{equation}
The pole at $s=1$ contributes $-\pi^2/(6t)$.  At $s=0$ there is a double pole.
Using
\[
 \zeta(0,x)=\frac12-x,
 \qquad
 \zeta'(0,x)=\log\Gamma(x)-\frac12\log(2\pi)
\]
gives the logarithmic and constant terms in~\eqref{eq:qx-infinite}.
At $s=-1$, the residue of $\Gamma$ and the values
$\zeta(0)=-1/2$, $\zeta(-1,x)=-B_2(x)/2$ yield $B_2(x)t/4$.
At $s=-2r$, $r\ge1$, use
\[
 \Res_{s=-2r}\Gamma(s)=\frac1{(2r)!},
 \quad
 \zeta(1-2r)=-\frac{B_{2r}}{2r},
 \quad
 \zeta(-2r,x)=-\frac{B_{2r+1}(x)}{2r+1}.
\]
The outer minus sign produces the coefficient in
\eqref{eq:qx-infinite}.  At negative odd integers below $-1$, a trivial zero of
$\zeta(s+1)$ cancels the gamma pole.  Shifting the contour gives the stated
sectorial remainder; the standard vertical estimates are uniform for $x$ on
compact subsets of the right half-plane.
\end{proof}

\begin{example}[The eta specialization]
At $x=1$, $B_{2r+1}(1)=0$ for $r\ge1$, so the algebraic expansion collapses to
\begin{equation}\label{eq:eta-algebraic}
 \log(q;q)_\infty
 \sim-\frac{\pi^2}{6t}+\frac12\log\frac{2\pi}{t}+\frac{t}{24}.
\end{equation}
There are no further algebraic powers.  The missing information is
exponentially small and is recovered exactly in \cref{chap:modular}.
\end{example}

\begin{example}[An integral shift]
For $x=2$,
\((q^2;q)_\infty=(q;q)_\infty/(1-q)
\).  Formula~\eqref{eq:qx-infinite} gives
\[
 \log(q^2;q)_\infty
 \sim-\frac{\pi^2}{6t}-\frac32\log t+\frac12\log(2\pi)
 +\frac{13}{24}t-\frac1{24}t^2+\cdots,
\]
which agrees term by term with subtracting
$\log(1-\e^{-t})$ from~\eqref{eq:eta-algebraic}.
\end{example}

\chapter{The \texorpdfstring{$q$}{q}-gamma, \texorpdfstring{$q$}{q}-beta, and \texorpdfstring{$q$}{q}-polygamma functions}\label{chap:qgamma-exp}

\section{All-order expansion of \texorpdfstring{$\Gamma_q$}{q-gamma}}

For $0<q<1$,
\[
 \Gamma_q(x)=(1-q)^{1-x}\frac{(q;q)_\infty}{(q^x;q)_\infty}.
\]
Combining \cref{thm:qx-infinite} with
\[
 \log\frac{1-\e^{-t}}t
 =-\frac t2+\sum_{r\ge1}\frac{B_{2r}}{2r(2r)!}t^{2r}
\]
gives the following explicit expansion.

\begin{theorem}[All-order $q\to1$ expansion of $\Gamma_q$]
\label{thm:qgamma-all}
For $\Re x>0$ and $q=\e^{-t}$,
\begin{align}
 \log\Gamma_{\e^{-t}}(x)
 \sim{}&\log\Gamma(x)+\frac{(x-1)(2-x)}4t
 \notag\\
 &+\sum_{r\ge1}\frac{B_{2r}}{2r(2r)!}
 \left(1-x+\frac{B_{2r+1}(x)}{2r+1}\right)t^{2r}.
 \label{eq:qgamma-all}
\end{align}
The expansion is locally uniform on compact subsets of $\Re x>0$ and in closed
subsectors of $\abs{\arg t}<\pi/2$.
\end{theorem}

\begin{proof}
Subtract~\eqref{eq:qx-infinite} from its $x=1$ specialization and add
$(1-x)\log(1-\e^{-t})$.  The $t^{-1}$, $\log(2\pi)$, and logarithmic terms
cancel.  The linear coefficient simplifies to
\[
 \frac1{24}-\frac{B_2(x)}4-\frac{1-x}{2}
 =\frac{(x-1)(2-x)}4.
\]
For $r\ge1$, the $x=1$ Bernoulli-polynomial term vanishes, and the two remaining
contributions combine exactly as displayed.
\end{proof}

The first terms are
\begin{align}
 \log\Gamma_q(x)
 ={}&\log\Gamma(x)+\frac{(x-1)(2-x)}4t
 \notag\\
 &+\frac1{24}\left(1-x+\frac{B_3(x)}3\right)t^2
 -\frac1{2880}\left(1-x+\frac{B_5(x)}5\right)t^4
 +\bigO(t^6).                                      \label{eq:qgamma-first}
\end{align}
For $x=3$, this reduces to
\[
 \log\Gamma_q(3)=\log(1+\e^{-t})
 =\log2-\frac t2+\frac{t^2}{8}-\frac{t^4}{192}+\cdots,
\]
providing an exact consistency check.

\section{\texorpdfstring{$q$}{q}-digamma and higher logarithmic derivatives}

Termwise differentiation in $x$ is legitimate on compact subsets of the right
half-plane.

\begin{corollary}[Expansion of the $q$-digamma function]
Let $\psi_q(x)=\partial_x\log\Gamma_q(x)$.  Then
\begin{align}
 \psi_{\e^{-t}}(x)
 \sim{}&\psi(x)+\frac{3-2x}{4}t
 \notag\\
 &+\sum_{r\ge1}\frac{B_{2r}}{2r(2r)!}
 \bigl(B_{2r}(x)-1\bigr)t^{2r}.                    \label{eq:qdigamma-all}
\end{align}
Further $x$-derivatives follow from
$B_n'(x)=nB_{n-1}(x)$.
\end{corollary}

\section{The \texorpdfstring{$q$}{q}-beta function}

Since
\(
B_q(x,y)=\Gamma_q(x)\Gamma_q(y)/\Gamma_q(x+y)
\), define
\[
 c_1(x)=\frac{(x-1)(2-x)}4,
 \qquad
 c_{2r}(x)=\frac{B_{2r}}{2r(2r)!}
 \left(1-x+\frac{B_{2r+1}(x)}{2r+1}\right).
\]
Then
\begin{equation}\label{eq:qbeta-all}
 \log\frac{B_{\e^{-t}}(x,y)}{B(x,y)}
 \sim\bigl(c_1(x)+c_1(y)-c_1(x+y)\bigr)t
 +\sum_{r\ge1}
 \bigl(c_{2r}(x)+c_{2r}(y)-c_{2r}(x+y)\bigr)t^{2r}.
\end{equation}
The linear coefficient simplifies to
\begin{equation}\label{eq:qbeta-linear}
 c_1(x)+c_1(y)-c_1(x+y)=\frac{xy-1}{2}.
\end{equation}
Bell polynomials exponentiate~\eqref{eq:qbeta-all} to any desired order.

\chapter{Modular and exponentially improved expansions}\label{chap:modular}

\section{The eta product at \texorpdfstring{$q=1$}{q = 1}}

Dedekind eta modularity gives an exact completion of
\eqref{eq:eta-algebraic}:
\begin{theorem}[Exact modular completion]\label{thm:eta-exact}
For $t>0$,
\begin{equation}\label{eq:eta-exact}
 (\e^{-t};\e^{-t})_\infty
 =\sqrt{\frac{2\pi}{t}}
 \exp\left(-\frac{\pi^2}{6t}+\frac{t}{24}\right)
 (\e^{-4\pi^2/t};\e^{-4\pi^2/t})_\infty.
\end{equation}
In particular, truncating after the displayed elementary factor incurs a
relative error $\bigO(\e^{-4\pi^2/t})$.
\end{theorem}

\begin{proof}
Apply $\eta(-1/\tau)=\sqrt{-\ii\tau}\,\eta(\tau)$ with
$\tau=\ii t/(2\pi)$ and use
$\eta(\tau)=q^{1/24}(q;q)_\infty$.
\end{proof}

This exact identity explains why every algebraic coefficient beyond $t/24$
vanishes: the remainder is not zero, but beyond all orders.

\section{The eta product on the negative real ray}

Let $x=\e^{-t}$.  Separating even and odd indices gives
\begin{equation}\label{eq:negative-eta-dissection}
 (-x;-x)_\infty
 =\frac{(x^2;x^2)_\infty^3}
 {(x;x)_\infty(x^4;x^4)_\infty}.
\end{equation}
Applying~\eqref{eq:eta-exact} at $t,2t,4t$ yields another exact exponentially
improved formula.

\begin{theorem}[Exact completion at $q=-1$]\label{thm:negative-eta}
For $t>0$, set
\(
Q_c=\exp(-c\pi^2/t)
\).  Then
\begin{equation}\label{eq:negative-eta-exact}
 (-\e^{-t};-\e^{-t})_\infty
 =\sqrt{\frac{\pi}{t}}
 \exp\left(-\frac{\pi^2}{24t}+\frac{t}{24}\right)
 \frac{(Q_2;Q_2)_\infty^3}{(Q_4;Q_4)_\infty(Q_1;Q_1)_\infty}.
\end{equation}
Consequently,
\begin{equation}\label{eq:negative-eta-leading}
 (-\e^{-t};-\e^{-t})_\infty
 =\sqrt{\frac{\pi}{t}}
 \exp\left(-\frac{\pi^2}{24t}+\frac{t}{24}\right)
 \left(1+\bigO(\e^{-\pi^2/t})\right).
\end{equation}
\end{theorem}

\begin{proof}
Identity~\eqref{eq:negative-eta-dissection} follows from
\[
 (-x;x^2)_\infty
 =\frac{(x^2;x^4)_\infty}{(x;x^2)_\infty},
 \quad
 (x;x^2)_\infty=\frac{(x;x)_\infty}{(x^2;x^2)_\infty},
 \quad
 (x^2;x^4)_\infty=\frac{(x^2;x^2)_\infty}{(x^4;x^4)_\infty}.
\]
Substitute the three modular completions and simplify the elementary factors.
\end{proof}

\section{An arbitrary eta cusp}

Let $h,m$ be coprime, $m>0$, and
\[
 \tau=\frac hm+\frac{\ii t}{2\pi},
 \qquad q=\e^{2\pi\ii\tau}=\e^{2\pi\ii h/m}\e^{-t}.
\]
Choose integers $a,b$ such that
\(
-ah-bm=1
\), and let
\(
\gamma=\begin{psmallmatrix}a&b\\m&-h\end{psmallmatrix}
\in\mathrm{SL}_2(\Z)
\).  Dedekind eta transformation gives the exact identity
\begin{equation}\label{eq:general-cusp-exact}
 (q;q)_\infty
 =\e^{-\pi\ii\tau/12}\,
 \epsilon(\gamma)^{-1}(m\tau-h)^{-1/2}
 \e^{\pi\ii\gamma\tau/12}
 \bigl(\e^{2\pi\ii\gamma\tau};
       \e^{2\pi\ii\gamma\tau}\bigr)_\infty,
\end{equation}
where $\epsilon(\gamma)$ is the eta multiplier and a consistent square-root
branch is understood.  Since
\[
 \gamma\tau=\frac am+\frac{2\pi\ii}{m^2t}.
\]
After the chosen Bezout relation is used, the final product differs from one by
$\bigO(\e^{-4\pi^2/(m^2t)})$.  In particular,
\begin{equation}\label{eq:general-cusp-magnitude}
 \abs{(q;q)_\infty}
 \sim\sqrt{\frac{2\pi}{mt}}
 \exp\left(-\frac{\pi^2}{6m^2t}+\frac{t}{24}\right).
\end{equation}
The cases $m=1$ and $m=2$ reproduce
\eqref{eq:eta-exact} and~\eqref{eq:negative-eta-leading}.

\section{A theta transseries by Poisson summation}

For
\[
 \Theta(z;q)=\sum_{n\in\Z}q^{n(n-1)/2}z^n,
 \qquad z=\e^u,
\]
completion of the square and Poisson summation give, for $t>0$,
\begin{equation}\label{eq:theta-poisson}
 \Theta(\e^u;\e^{-t})
 =\sqrt{\frac{2\pi}{t}}
 \exp\left(\frac{(u+t/2)^2}{2t}\right)
 \sum_{k\in\Z}
 \exp\left(-\frac{2\pi^2k^2}{t}
 +\frac{2\pi\ii k(u+t/2)}{t}\right).
\end{equation}
This is an exact transseries rather than a merely formal power expansion.  Its
dominant exponential changes when $u$ crosses Stokes lines, and multiplication
by the Jacobi triple product converts the statement into coupled asymptotics of
three infinite $q$-Pochhammer factors.

\chapter{Radial asymptotics at arbitrary roots of unity}\label{chap:root-radial}

\section{A fixed nonresonant parameter}

Let $\zeta$ be a primitive $m$th root of unity and approach it radially by
\[
 q=\zeta\e^{-t},
 \qquad t\to0,
 \qquad \Re t>0.
\]
For fixed $\abs a<1$, the Lambert series can be separated into indices divisible
by $m$ and indices not divisible by $m$.  The former create the singular
dilogarithmic scale; the latter are regular but remember the phase $\zeta$.

\begin{definition}[Cyclic quantum dilogarithm]
Set
\begin{equation}\label{eq:cyclic-dilog}
 D_\zeta(a)=\prod_{s=1}^{m-1}(1-\zeta^s a)^s.
\end{equation}
For $m=1$ the empty product is one.
\end{definition}

Define rational functions
\begin{equation}\label{eq:eulerian-rational}
 \mathcal E_j(u)=
 \left(u\frac{\dd}{\dd u}\right)^j\frac1{1-u}.
\end{equation}
Thus
\[
 \mathcal E_0(u)=\frac1{1-u},
 \quad
 \mathcal E_1(u)=\frac{u}{(1-u)^2},
 \quad
 \mathcal E_2(u)=\frac{u(1+u)}{(1-u)^3},
\]
and $\mathcal E_j(u)=\Li_{-j}(u)$ for $j\ge1$, interpreted as a rational
function when $u\ne1$.

\begin{theorem}[All-order fixed-$a$ root-of-unity expansion]
\label{thm:root-fixed-a}
Let $\abs a\le\rho<1$ and let $\zeta$ be primitive of order $m$.  Then, uniformly
for $a$ in that disk and for $t$ in a closed subsector of $\Re t>0$,
\begin{align}
 \log\poch a{\zeta\e^{-t}}\infty
 \sim{}&-\frac{\Li_2(a^m)}{m^2t}
 +\frac12\log(1-a^m)-\frac1m\log D_\zeta(a)
 \notag\\
 &+\sum_{j\ge1}C_j(a,\zeta)t^j,                  \label{eq:root-fixed-a}
\end{align}
where
\begin{align}
 C_j(a,\zeta)
 ={}&\frac{(-1)^{j+1}}{j!}
 \sum_{\substack{r\ge1\\m\nmid r}}
 a^r r^{j-1}\mathcal E_j(\zeta^r)
 \notag\\
 &-\frac{B_{j+1}m^{j-1}}{(j+1)!}
 \Li_{1-j}(a^m).                                  \label{eq:root-Cj}
\end{align}
All coefficient sums converge absolutely.  For $m=1$, the first line involving
$D_\zeta$ is empty and~\eqref{eq:root-fixed-a} reduces to
\cref{thm:infinite-fixed-a}.
\end{theorem}

\begin{proof}
Start from
\[
 L(t)=-\sum_{r\ge1}\frac{a^r}{r(1-\zeta^r\e^{-rt})}.
\]
When $r=m\ell$, expand
\[
 \frac1{1-\e^{-m\ell t}}
 =\frac1{m\ell t}+\frac12+
 \sum_{j\ge1}\frac{B_{j+1}}{(j+1)!}(m\ell t)^j.
\]
The singular, constant, and $t^j$ contributions are respectively
\[
 -\frac{\Li_2(a^m)}{m^2t},
 \qquad
 \frac1{2m}\log(1-a^m),
 \qquad
 -\frac{B_{j+1}m^{j-1}}{(j+1)!}\Li_{1-j}(a^m)t^j.
\]
For $m\nmid r$, the dilation identity
\[
 \frac1{1-u\e^{-x}}
 =\sum_{j\ge0}\frac{(-x)^j}{j!}\mathcal E_j(u)
\]
gives the first sum in~\eqref{eq:root-Cj}.  Its constant term is
\[
 -\sum_{\substack{r\ge1\\m\nmid r}}
 \frac{a^r}{r(1-\zeta^r)}.
\]
To simplify it, observe that for $m\nmid r$,
\[
 \sum_{s=1}^{m-1}s\zeta^{sr}=-\frac{m}{1-\zeta^r},
\]
while for $m\mid r$ the sum is $m(m-1)/2$.  Expanding
$\log D_\zeta(a)$ from~\eqref{eq:cyclic-dilog} therefore yields
\[
 -\sum_{\substack{r\ge1\\m\nmid r}}
 \frac{a^r}{r(1-\zeta^r)}
 =-\frac1m\log D_\zeta(a)
 +\frac{m-1}{2m}\log(1-a^m).
\]
Adding the constant from the divisible indices gives the constant in
\eqref{eq:root-fixed-a}.

For the remainder estimate, choose $\delta>0$ smaller than the distance from
$0$ to every singularity of
$x\mapsto(1-u\e^{-x})^{-1}$ with
$u\in\{\zeta,\ldots,\zeta^{m-1}\}$, and also choose $\delta<2\pi$ for the
resonant Bernoulli expansion.  Split the Lambert sum at
$r\le\delta/\abs t$.  On this first range the relevant Taylor remainders are
uniformly bounded by a polynomial in $r$ times the requested power of $t$.
On the complementary range, $\rho^r=\bigO(\exp(-c/\abs t))$ for some $c>0$,
so the tail is smaller than every algebraic power of $t$.  This is the same
small-range/exponentially-small-tail argument as in
\cref{thm:infinite-fixed-a}.
\end{proof}

\begin{remark}[Normalization used in the roots-of-unity literature]
If one writes $q=\zeta\e^{-\eps/m}$ instead, then $t=\eps/m$ and the leading
term becomes $-\Li_2(a^m)/(m\eps)$.  This is the normalization commonly used
in radial asymptotics of Nahm sums.  The cyclic factor
$D_\zeta(a)^{-1/m}$ is exactly the finite root-of-unity datum that survives
after the singular dilogarithm is removed; compare
Garoufalidis--Zagier~\cite{GaroufalidisZagier}.
\end{remark}

\section{The first coefficients}

Formula~\eqref{eq:root-Cj} is already computationally explicit.  For example,
\begin{align*}
 C_1={}&
 \sum_{m\nmid r}a^r\frac{\zeta^r}{(1-\zeta^r)^2}
 -\frac1{12}\Li_0(a^m),\\
 C_2={}&-\frac12
 \sum_{m\nmid r}a^r r
 \frac{\zeta^r(1+\zeta^r)}{(1-\zeta^r)^3},
\end{align*}
where the $B_3$ contribution vanishes.  The sums may be reduced to finite
linear combinations of $\Li_{1-j}(a\zeta^s)$ by discrete Fourier analysis.
Keeping the residue-class sums is often numerically better conditioned and
makes the distinction between resonant and nonresonant indices transparent.

\section{Coalescence and why a second scaling is necessary}

The hypothesis $\abs a\le\rho<1$ excludes $a\to\zeta^{-r}$, where one of the
finite cyclic factors vanishes.  It also excludes the eta specialization
$a=q$, for which $a^m\to1$.  In those regimes the dilogarithm in
\eqref{eq:root-fixed-a} develops a logarithmic singularity that must be combined
with the cyclic factor before expansion.  The resulting scale contains powers
of $t$, $\log t$, Bernoulli polynomials in the shift parameter, and an
exponentially small completion.  General all-order shifted formulae of this
kind underlie the root-of-unity analysis of Nahm sums
\cite{GaroufalidisZagier}; the eta specialization is treated exactly by
\eqref{eq:general-cusp-exact}.

A central open problem, formulated more sharply in
\cref{chap:frontiers}, is to construct one uniform normal form that contains
both the fixed-$a$ expansion~\eqref{eq:root-fixed-a} and the coalescing
$a^m\to1$ expansion without separate case distinctions.

\chapter{Other \texorpdfstring{$q$}{q}-analogues near \texorpdfstring{$q=1$}{q = 1}}\label{chap:other-q}

\section{The two standard \texorpdfstring{$q$}{q}-exponentials}

Use the normalization in the source monograph,
\[
 e_q(x)=\sum_{n\ge0}\frac{x^n}{[n]_q!}
 =\frac1{((1-q)x;q)_\infty},
 \qquad
 E_q(x)=\sum_{n\ge0}\frac{q^{\binom n2}x^n}{[n]_q!}
 =(-(1-q)x;q)_\infty.
\]
Let $h=1-q$.  For $0<q<1$ and $\abs{hx}<1$, with logarithms continued from
$x=0$, their logarithms have the exact convergent representations
\begin{align}
 \log e_q(x)&=\sum_{m\ge1}\frac{h^m x^m}{m(1-(1-h)^m)},\label{eq:log-eq-exact}\\
 \log E_q(x)&=\sum_{m\ge1}\frac{(-1)^{m+1}h^m x^m}
 {m(1-(1-h)^m)}.                                  \label{eq:log-Eq-exact}
\end{align}
At any prescribed order in $h$, only finitely many $m$ contribute.
Consequently these equations are also coefficient-generation algorithms.
Expanding gives
\begin{align}
 \log e_q(x)
 ={}&x+\frac{x^2}{4}h+
 \left(\frac{x^2}{8}+\frac{x^3}{9}\right)h^2
 +\frac{x^2(9x^2+16x+9)}{144}h^3+\bigO(h^4),
 \label{eq:eq-first}\\
 \log E_q(x)
 ={}&x-\frac{x^2}{4}h+
 \left(-\frac{x^2}{8}+\frac{x^3}{9}\right)h^2
 -\frac{x^2(9x^2-16x+9)}{144}h^3+\bigO(h^4).
 \label{eq:Eq-first}
\end{align}
Bell exponentiation converts these logarithmic formulae to expansions of the
functions themselves.  The relation $e_q(x)E_q(-x)=1$ is visible term by term.

\begin{proof}[Derivation of \eqref{eq:log-eq-exact}--\eqref{eq:log-Eq-exact}]
Apply the Lambert identity~\eqref{eq:lambert-log-poc} with
$a=\mp hx$ and $q=1-h$, remembering the minus sign for the reciprocal in
$e_q$.
\end{proof}

\section{The \texorpdfstring{$q$}{q}-difference operator as an operator-valued \texorpdfstring{$q$}{q}-number}

Let $\Theta_x=x\frac{\dd}{\dd x}$.  For $q=\e^{-t}$,
\begin{equation}\label{eq:qderivative-operator}
 D_qf(x)=\frac{f(x)-f(qx)}{(1-q)x}
 =x^{-1}\frac{1-\e^{-t\Theta_x}}{1-\e^{-t}}f(x)
 =x^{-1}[\Theta_x]_{\e^{-t}}f(x).
\end{equation}
Using~\eqref{eq:qnumber-qone} with the scalar $x$ replaced by the commuting
operator $\Theta_x$ gives an all-order expansion.  The first terms are
\begin{align}
 D_{\e^{-t}}
 =x^{-1}\Bigl[&\Theta_x
 -\frac{t}{2}\Theta_x(\Theta_x-1)
 +\frac{t^2}{12}\Theta_x(\Theta_x-1)(2\Theta_x-1)
 \notag\\
 &-\frac{t^3}{24}\Theta_x^2(\Theta_x-1)^2
 +\bigO(t^4)\Bigr].                                \label{eq:qderivative-first}
\end{align}
Since $x^{-1}\Theta_x=\dd/\dd x$, the first term is the classical derivative.
Formula~\eqref{eq:qderivative-first} is valid coefficientwise on every fixed
finite-degree polynomial, with a remainder depending on that degree, and on
analytic functions only under local convergence and operator-domain
hypotheses.  It is an asymptotic operator expansion, not a finite Taylor
identity in $t$.

\section{Jackson integration by Euler--Maclaurin}

For $q=\e^{-t}$,
\[
 \int_0^a f(x)\dd_qx
 =(1-\e^{-t})a\sum_{n\ge0}\e^{-nt}f(a\e^{-nt}).
\]
Put $g(s)=a\e^{-s}f(a\e^{-s})$.  Suppose that $f$ is sufficiently smooth on
$[0,a]$ and that the transformed derivatives of $g$ required below are
integrable and decay at infinity.  Euler--Maclaurin then gives
\begin{align}
 t\sum_{n\ge0}g(nt)
 \sim\int_0^\infty g(s)\dd s+\frac t2g(0)
 -\sum_{r\ge1}\frac{B_{2r}t^{2r}}{(2r)!}g^{(2r-1)}(0).
 \label{eq:jackson-em}
\end{align}
Multiplication by $(1-\e^{-t})/t$ yields
\begin{align}
 \int_0^a f(x)\dd_{\e^{-t}}x
 ={}&\int_0^a f(x)\dd x
 +\frac t2\left(af(a)-\int_0^a f(x)\dd x\right)
 \notag\\
 &+t^2\left(
 \frac{a^2f'(a)}{12}-\frac{af(a)}6
 +\frac16\int_0^a f(x)\dd x\right)+\bigO(t^3).
 \label{eq:jackson-first}
\end{align}
The full expansion is an explicit endpoint differential series.  It provides a
systematic bridge between Jackson $q$-integrals and ordinary quadrature, and is
particularly convenient for perturbing $q$-beta integrals.

\chapter{Basic hypergeometric series as differential-operator deformations}
\label{chap:hypergeom-deform}

\section{A coefficientwise all-order operator}

Consider the balanced-length family
\begin{equation}\label{eq:rphi}
 F_q(z)={}_r\phi_{r-1}
 \left[\begin{matrix}q^{\alpha_1},\ldots,q^{\alpha_r}\\
 q^{\beta_1},\ldots,q^{\beta_{r-1}}
 \end{matrix};q,z\right],
\end{equation}
with $\beta_r=1$ added to account for the factor $(q;q)_n$.  Its classical
limit is
\[
 F(z)={}_rF_{r-1}
 \left[\begin{matrix}\alpha_1,\ldots,\alpha_r\\
 \beta_1,\ldots,\beta_{r-1}
 \end{matrix};z\right].
\]
Let $\theta=z\frac{\dd}{\dd z}$ and define the polynomial
\begin{align}
 P_{2s}(X)=\sum_{i=1}^{r}
 \frac{B_{2s+1}(\alpha_i+X)-B_{2s+1}(\alpha_i)}{2s+1}
 \notag\\
 -\sum_{j=1}^{r}
 \frac{B_{2s+1}(\beta_j+X)-B_{2s+1}(\beta_j)}{2s+1}.
 \label{eq:P-hypergeom}
\end{align}
Because the two parameter lists have equal length, the leading $X^{2s+1}$
terms cancel, so $P_{2s}$ has degree at most $2s$.

\begin{theorem}[Formal differential-operator expansion]
\label{thm:hypergeom-operator}
Let
\[
 A=\sum_{i=1}^{r}\alpha_i-\sum_{j=1}^{r}\beta_j.
\]
Assume that the denominator parameters avoid nonpositive integers and that
$q$ approaches one through a neighborhood in which all denominator shifted
factorials occurring in the coefficients under consideration are nonzero.
Then, coefficientwise in $z$ and asymptotically as $q=\e^{-t}\to1$,
\begin{equation}\label{eq:hypergeom-operator}
 F_q(z)\sim
 \exp\left(
 -\frac{A}{2}t\theta+
 \sum_{s\ge1}\frac{B_{2s}}{2s(2s)!}
 P_{2s}(\theta)t^{2s}
 \right)F(z).
\end{equation}
On every compact disk on which the defining series and a sufficient number of
its parameterwise majorants converge uniformly, the expansion may be
differentiated and summed termwise to any fixed order.
\end{theorem}

\begin{proof}
The coefficient of $z^n$ in~\eqref{eq:rphi} is a product of ratios
$(q^\alpha;q)_n/(q^\beta;q)_n$.  Apply
\cref{thm:scaled-finite}.  The linear difference is $An$, and the $2s$th power
sum difference is exactly $P_{2s}(n)$.  Since $\theta z^n=nz^n$, replacing $n$
by $\theta$ packages the coefficientwise multiplier into the operator
exponential~\eqref{eq:hypergeom-operator}.
\end{proof}

\section{The first two perturbative orders}

Put
\[
 B=\sum_{i=1}^{r}\alpha_i^2-
   \sum_{j=1}^{r}\beta_j^2.
\]
Since $P_2(n)=An(n-1)+Bn$, expansion of the operator exponential gives
\begin{align}
 F_q(z)
 ={}&F(z)-\frac A2t\,\theta F(z)
 \notag\\
 &+\frac{t^2}{24}
 \left((3A^2+A)\theta^2+(B-A)\theta\right)F(z)
 +\bigO(t^3).                                      \label{eq:hypergeom-second}
\end{align}
If $A=0$, the entire first-order correction vanishes.  If also $B=0$, the
expansion starts at order $t^4$ in the logarithmic operator and at order $t^4$
in the function itself.  These cancellation criteria offer a quick way to
identify especially faithful $q$-deformations of classical hypergeometric
functions.

\begin{example}[$q$-binomial theorem]
For ${}_1\phi_0(q^\alpha;-;q,z)$, one has
$F(z)=(1-z)^{-\alpha}$,
$A=\alpha-1$, and $B=\alpha^2-1$.  Formula
\eqref{eq:hypergeom-second} gives the first corrections entirely in terms of
$\theta(1-z)^{-\alpha}$ and $\theta^2(1-z)^{-\alpha}$, without re-expanding the
infinite product quotient separately.
\end{example}

\section{Terminating series and roots of unity}

For a terminating basic hypergeometric series, each summand is a rational
function of $q$ and the parameters.  At a chosen root of unity, cyclotomic
factorization decides whether apparent zero-over-zero factors cancel, whether
the summand has a genuine zero, or whether it has a genuine pole.  The reliable
procedure is therefore:

\begin{enumerate}[label=\arabic*.]
\item factor every finite shifted factorial into cyclotomic factors in the
      numerator and denominator;
\item compare the local valuations and extract the resulting power of
      $q-\zeta$ (negative powers signal an actual pole rather than a removable
      singularity);
\item evaluate the regular nonzero quotient in $\Q(\zeta)$;
\item use Bell polynomials for its higher jet and then sum the finitely many
      terms whenever the complete terminating expression is regular.
\end{enumerate}

Direct numerical substitution into an uncancelled quotient can lose every
significant digit.  The cyclotomic method computes the exact local algebra
first and only then performs numerical evaluation.

\chapter{A double-scaling expansion for Gaussian coefficients}
\label{chap:double-scaling}

\section{Why fixed \texorpdfstring{$n$}{n} is not uniform}

The convergent fixed-$n$ expansion~\eqref{eq:gaussian-log-all} has radius
$2\pi/n$ in $t$.  It is therefore not uniform when $n\to\infty$.  The natural
transition coordinate is
\[
 q=\e^{-\tau/N},
 \qquad k=\alpha N,
 \qquad 0<\alpha<1,
\]
with $\tau$ fixed.  This regime interpolates between the classical binomial
coefficient ($\tau\to0$) and a genuinely dilogarithmic $q$-deformation.

Define
\begin{equation}\label{eq:I-tau}
 I_\tau(\beta)=\int_0^\beta\log(1-\e^{-\tau x})\dd x
 =\frac{\Li_2(\e^{-\beta\tau})-\pi^2/6}{\tau}.
\end{equation}

\section{A finite \texorpdfstring{$q$}{q}-factorial with a moving endpoint}

\begin{theorem}[All-order moving-endpoint expansion]
\label{thm:moving-qfactorial}
Fix $\tau>0$ and $\beta>0$, and let $\beta N$ be integral.  With
$q=\e^{-\tau/N}$,
\begin{align}
 \log(q;q)_{\beta N}
 \sim{}&N I_\tau(\beta)
 +\frac12\log\left(\frac{2\pi N(1-\e^{-\beta\tau})}{\tau}\right)
 +\frac{\tau}{24N}\cothh\left(\frac{\beta\tau}{2}\right)
 \notag\\
 &+\sum_{r\ge2}
 \frac{B_{2r}\tau^{2r-1}}{(2r)!N^{2r-1}}
 \Li_{2-2r}(\e^{-\beta\tau}).                    \label{eq:moving-qfactorial}
\end{align}
The expansion is uniform when $(\beta,\tau)$ ranges over a compact subset of
$(0,\infty)^2$.  After the $N^{-1}$ term, only odd powers of $N^{-1}$ occur.
\end{theorem}

\begin{proof}
Write
\[
 \log(1-\e^{-\tau x})=\log(\tau x)+g(\tau x),
 \qquad
 g(u)=\log\frac{1-\e^{-u}}u.
\]
Then
\begin{align*}
 \log(q;q)_{\beta N}
 &=\beta N\log\frac{\tau}{N}+\log(\beta N)!
   +\sum_{j=1}^{\beta N}g\left(\frac{\tau j}{N}\right).
\end{align*}
Apply Stirling's expansion to the factorial and Euler--Maclaurin to the smooth
last sum.  Their leading terms combine into $NI_\tau(\beta)$.  The endpoint
terms simplify to the logarithm in~\eqref{eq:moving-qfactorial}.

For completeness, let $H(x)=g(\tau x)$.  The coefficient of
$N^{-(2r-1)}$ before simplification is
\begin{equation}\label{eq:moving-unsimplified}
 \frac{B_{2r}}{(2r)!}
 \left(
 \frac{(2r-2)!}{\beta^{2r-1}}
 +H^{(2r-1)}(\beta)-H^{(2r-1)}(0)
 \right).
\end{equation}
For $r=1$, $H'(0)=-\tau/2$, and
\eqref{eq:moving-unsimplified} becomes
$\tau\cothh(\beta\tau/2)/24$.  For $r\ge2$,
$H^{(2r-1)}(0)=0$ and
\[
 H^{(2r-1)}(\beta)
 =\tau^{2r-1}\Li_{2-2r}(\e^{-\beta\tau})
 -\frac{(2r-2)!}{\beta^{2r-1}}.
\]
The rational endpoint term cancels the Stirling term, leaving exactly the
polylogarithm in~\eqref{eq:moving-qfactorial}.  Standard uniform
Euler--Maclaurin bounds prove the stated uniformity.
\end{proof}

\section{Gaussian coefficients}

Define the deformed entropy
\begin{align}
 \Phi(\alpha,\tau)
 &\defeq I_\tau(1)-I_\tau(\alpha)-I_\tau(1-\alpha)\notag\\
 &=\frac{\Li_2(\e^{-\tau})-\Li_2(\e^{-\alpha\tau})
 -\Li_2(\e^{-(1-\alpha)\tau})+\pi^2/6}{\tau}.
 \label{eq:deformed-entropy}
\end{align}
Also put
\begin{align}
 A_1(\beta;\tau)&=\frac{\tau}{24}
 \cothh\left(\frac{\beta\tau}{2}\right),\label{eq:A1-double}\\
 A_{2r-1}(\beta;\tau)&=
 \frac{B_{2r}\tau^{2r-1}}{(2r)!}
 \Li_{2-2r}(\e^{-\beta\tau})\qquad(r\ge2).
 \label{eq:Ar-double}
\end{align}

\begin{corollary}[All-order double scaling of $\qbinomq Nk$]
\label{cor:double-gaussian}
Let $k=\alpha N$ be integral, with $(\alpha,\tau)$ in a compact subset of
$(0,1)\times(0,\infty)$.  Then
\begin{align}
 \log\qbinombase N{\alpha N}{\e^{-\tau/N}}
 \sim{}&N\Phi(\alpha,\tau)
 -\frac12\log\frac{2\pi N}{\tau}
 \notag\\
 &+\frac12\log
 \frac{1-\e^{-\tau}}
 {(1-\e^{-\alpha\tau})(1-\e^{-(1-\alpha)\tau})}
 \notag\\
 &+\sum_{r\ge1}\frac{
 A_{2r-1}(1;\tau)-A_{2r-1}(\alpha;\tau)
 -A_{2r-1}(1-\alpha;\tau)}{N^{2r-1}}.
 \label{eq:double-gaussian}
\end{align}
\end{corollary}

\begin{proof}
Subtract the two denominator instances of
\eqref{eq:moving-qfactorial} from the numerator instance.
\end{proof}

The first correction is therefore
\begin{equation}\label{eq:double-first-correction}
 \frac{\tau}{24N}
 \left[
 \cothh\frac\tau2-
 \cothh\frac{\alpha\tau}{2}-
 \cothh\frac{(1-\alpha)\tau}{2}
 \right].
\end{equation}
The next is
\begin{equation}\label{eq:double-third-correction}
 -\frac{\tau^3}{720N^3}
 \left[
 \Li_{-2}(\e^{-\tau})-
 \Li_{-2}(\e^{-\alpha\tau})-
 \Li_{-2}(\e^{-(1-\alpha)\tau})
 \right].
\end{equation}

\section{Recovery of ordinary Stirling asymptotics}

As $\tau\to0$,
\[
 \Phi(\alpha,\tau)\longrightarrow
 -\alpha\log\alpha-(1-\alpha)\log(1-\alpha).
\]
The logarithmic prefactor in~\eqref{eq:double-gaussian} tends to
\[
 -\frac12\log\bigl(2\pi N\alpha(1-\alpha)\bigr),
\]
and~\eqref{eq:double-first-correction} tends to
\[
 \frac1{12N}
 \left(1-\frac1\alpha-\frac1{1-\alpha}\right).
\]
These are exactly the entropy, Gaussian prefactor, and first Stirling
correction for $\binom N{\alpha N}$.  Hence the double-scaling formula is a
uniform $q$-deformation of classical binomial asymptotics rather than a separate
limit formula.

\section{Multinomial double scaling}

Let $\alpha_1+\cdots+\alpha_s=1$ with every $\alpha_i>0$ and
$n_i=\alpha_iN$ integral.  The same subtraction gives
\begin{align}
 \log\qmultinom{N}{\alpha_1N,\ldots,\alpha_sN}
 \bigg|_{q=\e^{-\tau/N}}
 \sim{}&N\left(I_\tau(1)-\sum_{i=1}^{s}I_\tau(\alpha_i)\right)
 \notag\\
 &+\frac{1-s}{2}\log\frac{2\pi N}{\tau}
 +\frac12\log\frac{1-\e^{-\tau}}
 {\prod_i(1-\e^{-\alpha_i\tau})}
 \notag\\
 &+\sum_{r\ge1}\frac{
 A_{2r-1}(1;\tau)-\sum_iA_{2r-1}(\alpha_i;\tau)}
 {N^{2r-1}}.                                      \label{eq:double-multinomial}
\end{align}
The formula is uniform when the composition remains in a compact subset of the
open simplex.

\chapter{Implications for the ProveIt and Fabius-function programme}
\label{chap:fabius-implications}

The primary purpose of this report is the $q$-series expansion problem itself,
but several outputs are directly reusable in the surrounding Fabius-function
research.

\section{Effective coefficient generation}

Every finite expansion in the report reduces to operations already friendly to
formalization and exact computation:

\begin{enumerate}[label=\textup{(\alph*)}]
\item finite power sums and Bernoulli polynomials;
\item cyclotomic factorization over $\Z[q]$;
\item logarithmic derivatives of nonvanishing polynomials;
\item complete Bell polynomials or the equivalent triangular recurrence;
\item exact rational or cyclotomic-number arithmetic.
\end{enumerate}

Thus local sensitivity of a formula involving $q$-binomial coefficients at the
dyadic points $q=1/2$ or $q=2$ can be computed by
\cref{thm:regular-bell,prop:gaussian-regular-jet}, while a deformation of the
base toward one can use the much smaller Bernoulli data in
\cref{thm:gaussian-complete}.  No numerical differentiation is required.

\section{Reciprocal dyadic normalizations}

The Fabius and Rvachev formulae frequently contain powers of two together with
Gaussian coefficients or products at bases $2$ and $1/2$.  Identities
\eqref{eq:poc-reciprocity} and~\eqref{eq:gaussian-reciprocity} make the
normalization explicit before numerical evaluation.  Moving a large power of
$q$ outside the polynomial leaves a convergent series in $q^{-1}$ and prevents
overflow or catastrophic cancellation.  The same transformation is useful in
formal proofs because it replaces a statement at $q>1$ by one inside the
convergence disk.

\section{Moment and cumulant organization}

Gaussian coefficients occurring in moment formulae can be expanded in a
centered form whose odd logarithmic coefficients vanish.  The exact cumulants
in \cref{cor:inversion-cumulants} provide a compact replacement for repeated
ad hoc differentiation.  In particular, when a Fabius-related expression
contains a weighted sum of Gaussian coefficients near $q=1$, its perturbative
coefficients can often be organized as finite combinations of inversion
cumulants and hence of the factored polynomials
\eqref{eq:delta2}--\eqref{eq:delta6}.

\section{Root-of-unity cancellation as exact algebra}

Alternating signs and Thue--Morse structures naturally bring $q=-1$ and other
roots of unity into intermediate formulae.  The cyclotomic-jet method separates
genuine zeros from removable factorial quotients.  The closed first jet in
\cref{thm:gaussian-minus-one-jet} is especially useful when a limiting formula
contains one vanishing Gaussian factor: the limit is then determined by a
single integer binomial coefficient rather than by a symbolic l'Hopital
calculation.

\chapter{Numerical validation and reproducible computation}\label{chap:numerics}

\section{What was checked}

The companion program \texttt{q\_expansion\_experiments.py} uses exact integer
polynomial arithmetic for finite Gaussian coefficients and 80-digit
\texttt{mpmath} arithmetic for infinite products.  It performs the following
checks.

\begin{enumerate}
\item It derives $\Delta_2,\Delta_4,\Delta_6$ symbolically and reduces symmetric
      expressions to polynomials in $n$ and $D=k(n-k)$.
\item It expands $\prod_{j\ge1}(1-aq^j)$ through a requested degree and verifies
      the weighted distinct-partition coefficients at $q=0$.
\item It constructs every $\qbinomq nk$ with $0\le k\le n\le20$ by Pascal
      recursion, checks the normalized $q=1$ expansion through fourth order,
      and checks both formulas in \cref{thm:gaussian-minus-one-jet}.  Each suite
      covers all 230 index pairs and passes exactly.
\item It evaluates the Mellin expansion~\eqref{eq:qx-infinite}, the
      $q$-gamma expansion~\eqref{eq:qgamma-all}, and the root-of-unity
      expansion~\eqref{eq:root-fixed-a} against directly summed logarithmic
      products.
\item It compares the double-scaling formula
      \eqref{eq:double-gaussian} through $N^{-1}$ and through $N^{-3}$.
\end{enumerate}

\section{Representative errors}

\begin{table}[htbp]
\centering
\caption{High-precision checks produced by the companion program.  Errors are
absolute errors in logarithms, except in the final two rows where they are
signed.}
\label{tab:numerical-errors}
\begin{tabular}{@{}lll@{}}
\toprule
Formula and parameters & Truncation & Error\\
\midrule
$(q^x;q)_\infty$, $x=2.5$, $t=0.1$ & through $t^6$ & $2.65\times10^{-14}$\\
$\Gamma_q(x)$, $x=3.7$, $t=0.08$ & through $t^6$ & $5.01\times10^{-13}$\\
root $m=3$, $a=0.2$, $t=0.05$ & through $t^6$ & $3.48\times10^{-11}$\\
$N=100$, $\alpha=0.4$, $\tau=1.7$ & through $N^{-1}$ & $+5.35\times10^{-8}$\\
 same double scaling & through $N^{-3}$ & $-8.69\times10^{-12}$\\
\bottomrule
\end{tabular}
\end{table}

The root-of-unity computation combines complex logarithm branches and a
truncated residue-class coefficient sum.  At the single displayed
double-scaling test point, adding the explicit $N^{-3}$ term reduces the error
by roughly four orders of magnitude.  These isolated evaluations do not test a
convergence rate, uniformity, or the order of the next remainder.

\section{Stable evaluation rules}

The experiments suggest the following implementation discipline.

\begin{enumerate}[label=\arabic*.]
\item Evaluate logarithms of long products and exponentiate only at the end.
\item Near $q=1$, use $t=-\log q$ computed by a compensated
      \texttt{log1p}; use \texttt{expm1} for $1-q^j$.
\item Near a root of unity, cancel cyclotomic zero factors symbolically before
      floating-point evaluation.
\item At $\abs q>1$, reverse the base and remove the explicit power of $q$
      before evaluation.
\item Choose the truncation order asymptotically: terms of a divergent
      Bernoulli series should be stopped near their least magnitude, not at a
      fixed arbitrarily high order.
\end{enumerate}

\chapter{Conjectures and research directions}\label{chap:frontiers}

The formulae proved in this TeX report expose several possible extensions.
The statements below are deliberately separated from the proved results and
have no Lean counterparts.

\section{A uniform coalescing root-of-unity normal form}

The fixed-$a$ theorem~\eqref{eq:root-fixed-a} fails uniformly as $a^m\to1$,
while shifted root-of-unity formulae reorganize the same object in a scale
containing $\log t$ and gamma factors.  The recent exact gamma-product formula
of Arabi Ardehali--Rosengren provides a new route to such uniformity near
$q=1$~\cite{ArdehaliRosengren}.

\begin{remark}[Informal uniform cyclic-gamma research target]\label{conj:uniform-root}
Let $\zeta$ be primitive of order $m$, let
$q=\zeta\e^{-t}$, and let
$a=a(t)$ satisfy
\(
1-a(t)^m=\bigO(t^\sigma)
\)
for a fixed $\sigma>0$.  One seeks a precise choice of branches, cut domain,
normalization, and remainder class such that, after extracting an explicit
product of
\begin{enumerate}[label=\textup{(\roman*)}]
\item $\exp(-\Li_2(a^m)/(m^2t))$,
\item a cyclic quantum dilogarithm,
\item finitely many ordinary gamma factors in the scaled variable
      $(1-a^m)/t$, and
\item an elementary power of $t$,
\end{enumerate}
there would be a single coefficient expansion, uniform for the scaled variable in
compact subsets of a cut plane, that reduces to
\eqref{eq:root-fixed-a} when $(1-a^m)/t\to\infty$ and to the shifted
root-of-unity expansion when $(1-a^m)/t=\bigO(1)$.
This is not yet a theorem-grade conjecture: until those choices are fixed, it
is a research template rather than a uniquely falsifiable statement.
\end{remark}

One possible derivation route is residue-class dissection plus an exact
Euler--Maclaurin formula with a singular endpoint retained rather than
expanded.  This would eliminate the current need to choose between ``fixed
parameter'' and ``coalescing parameter'' cases.

\section{Higher \texorpdfstring{$q$}{q}-Lucas jets}

The value at a root of unity is controlled by $q$-Lucas, and the zero order is a
single carry.  The theorem at $q=-1$ shows that the first derivative also
factorizes into a small polynomial in quotient digits times an ordinary
binomial coefficient.

\begin{conjecture}[Polynomial cyclotomic jets]\label{conj:polynomial-jets}
Fix $m$, a primitive $m$th root $\zeta$, residues
$0\le r,s<m$, and a jet order $J$.  Write $n=am+r$ and $k=bm+s$.  After
extracting the possible simple factor $q-\zeta$ and the natural leading
$q$-Lucas factor, the first $J$ normalized Taylor coefficients of
$\qbinomq nk$ at $q=\zeta$ are polynomials in $a$ and $b$ with coefficients in
$\Q(\zeta)$, of degree bounded linearly in $J$.
\end{conjecture}

Block the factorial products into intervals of length $m$.  Each block has a
Taylor logarithm polynomial in the block index, so Faulhaber summation strongly
suggests the conjecture.  A proof would give a genuine ``differential
$q$-Lucas theorem'' and an efficient algorithm independent of the size of
$n$.

\section{Rigorous complex double scaling}

\Cref{cor:double-gaussian} is proved for positive real $\tau$ with uniformity on
compact sets.  The same Euler--Maclaurin derivation should extend to complex
$\tau$ wherever the segments
$\{\beta\tau:0\le\beta\le1\}$ arising in the endpoint integrals avoid the
nonzero lattice $2\pi\ii\Z$.  Natural limiting cuts are therefore the imaginary
rays beginning at $\pm2\pi\ii$ (or at $\pm2\pi\ii/\beta$ for an individual
endpoint parameter $\beta$).  For finite $N$, zeros and poles occur at rational
points corresponding to roots of unity; as $N$ grows, these points become
dense along the limiting cuts.  A worthwhile next step is a uniform theorem on
compact subsets of the cut $\tau$-plane, followed by matched asymptotics as a
cut is approached.

The double-scaling action $\Phi(\alpha,\tau)$ is also a natural candidate for a
large-deviation rate function for inversion statistics under exponential
tilting.  Differentiation with respect to $\log q=-\tau/N$ gives limiting
cumulant densities, while Legendre transformation should produce an explicit
dilogarithmic rate function.

\section{Resurgence and exponentially small sectors}

Bernoulli coefficients grow factorially, so the fixed-$a$ asymptotic series is
normally divergent after analytic continuation beyond its local convergence
setting.  Exact eta modularity identifies the complete exponentially small
sector in one special case.  Recent work shows that suitable $q$-Pochhammer
combinations possess resurgent and quantum-modular structures
\cite{FantiniRella,ChengComanFantiniRella}.

Promising concrete tasks are:
\begin{enumerate}[label=\textup{(\alph*)}]
\item compute the Borel transform of the coefficient series
      in~\eqref{eq:root-fixed-a} for general $m$ and locate its singularities;
\item express the Stokes constants directly through residue-class Fourier
      transforms of $\mathcal E_j(\zeta^r)$;
\item determine when Dirichlet-character weighted cyclic products cancel all
      but one tower of Borel singularities;
\item compare the exact gamma-product completion
      \cite{ArdehaliRosengren} with median Borel summation.
\end{enumerate}

\section{Automatic proof-producing expansion}

The finite theory suggests a concrete Lean formalization project: a future
verified expansion compiler would use the following pipeline:
\[
 \begin{aligned}
 \text{product or cyclotomic factorization}
 &\longrightarrow \text{logarithmic jets}
 \longrightarrow \text{Bell recurrence}\\
 &\longrightarrow \text{Taylor certificate}.
 \end{aligned}
\]
For Gaussian coefficients at $q=1$, Bernoulli power sums replace polynomial
factorization; at roots of unity, exact arithmetic takes place in a cyclotomic
extension.  The certificate can be an identity in a truncated power-series
ring, so no analytic infrastructure is needed for finite products.  Analytic
remainder bounds can then be layered on separately.

\section{Optimal truncation and certified error bounds}

For numerical work, an asymptotic formula is most valuable when paired with a
computable remainder estimate.  The fixed-$a$ Lambert proof can be made to
give a simple geometric bound, but the report does not derive a sharp explicit
constant at high order.  A useful
research target is an adaptive bound that combines:
\begin{enumerate}[label=\textup{(\roman*)}]
\item the distance to the nearest singularity in the $t$-plane;
\item exact Bernoulli-number bounds;
\item the magnitude of the relevant polylogarithm;
\item an exponentially small tail estimated through a modular or gamma-product
      transform.
\end{enumerate}
Such a bound would automatically select the least term and produce a certified
number of correct digits.

\chapter{Conclusions}

The expansions of $q$-analogues are governed by a small number of reusable
mechanisms.

First, finite products are algebraic.  Logarithmic derivatives plus Bell
polynomials solve every regular local problem, while factor extraction solves
every resonant one.  At $q=0$, the same products become weighted partition
generators; at infinity, reciprocity reverses those coefficients.  Cyclotomic
factorization is the correct coordinate system at roots of unity.

Second, the coordinate $q=\e^{-t}$ linearizes the multiplicative geometry near
$q=1$.  Bernoulli numbers enter through the universal factor
$1-\e^{-x}$, power sums record finite endpoints, and polylogarithms record
infinite Lambert sums.  Bell polynomials then convert logarithmic data into
ordinary coefficients.  This single scheme covers finite shifted factorials,
Gaussian and multinomial coefficients, $q$-gamma and $q$-beta functions, and
basic hypergeometric series.

Third, modularity and roots of unity require information beyond algebraic
powers.  Eta transformations expose exact exponentially small corrections,
while residue-class splitting produces cyclic quantum dilogarithms.  The
boundary between fixed and coalescing parameters is the main remaining
uniformity problem.

Finally, the scale $q=\e^{-\tau/N}$ shows how fixed-degree perturbation theory
merges with large combinatorial systems.  Its all-order odd inverse-power
expansion is both a $q$-analogue of Stirling's formula and a dilogarithmic
large-deviation expansion.  Together with the exact cyclotomic jets, it offers
a concrete path from symbolic identities to certified asymptotic computation
and formal proof.


\appendix
\renewcommand{\theHsection}{appendix.\Alph{section}}
\renewcommand{\theHsubsection}{\theHsection.\arabic{subsection}}
\renewcommand{\theHtheorem}{\theHsection.\arabic{theorem}}

\chapter{Coefficient extraction and conversion tables}\label{app:coefficients}

\section{From a logarithmic expansion to an ordinary expansion}

The most economical way to calculate high-order coefficients is to retain the
logarithm until the final step.  The following recurrence is equivalent to the
complete exponential Bell-polynomial formula, but it is usually faster in a
computer algebra system and simpler to formalize.

\begin{proposition}[Logarithm-to-product recurrence]\label{prop:log-product-recurrence}
Suppose
\[
 L(z)=\sum_{r\ge1}\ell_r z^r,
 \qquad
 \exp L(z)=\sum_{m\ge0}c_mz^m.
\]
Then $c_0=1$ and, for $m\ge1$,
\begin{equation}\label{eq:log-product-recurrence}
 c_m=\frac1m\sum_{r=1}^{m}r\ell_r c_{m-r}.
\end{equation}
If instead
$L(z)=\sum_{r\ge1}\kappa_rz^r/r!$, then
\[
 c_m=\frac1{m!}\Bell_m(\kappa_1,\ldots,\kappa_m),
\]
where $\Bell_m$ is the complete exponential Bell polynomial.
\end{proposition}

\begin{proof}
Differentiate $C(z)=\exp L(z)$ to obtain $C'=L'C$.  Equating the coefficient
of $z^{m-1}$ gives
$m c_m=\sum_{r=1}^{m}r\ell_r c_{m-r}$.  The Bell-polynomial statement is the
standard exponential generating-function form of the same identity.
\end{proof}

When the original function has a zero of order $\nu$, first write
\[
 F(z)=z^\nu C\exp L(z),\qquad C\ne0,
\]
and apply the recurrence to the unit factor.  This one preliminary extraction
is what makes the same algorithm valid at regular and resonant points.

\section{Conversion between \texorpdfstring{$t=-\log q$}{t = -log q} and \texorpdfstring{$h=q-1$}{h = q - 1}}

Let $s(m,r)$ denote the signed Stirling number of the first kind, determined by
\[
 x(x-1)\cdots(x-m+1)=\sum_{r=0}^{m}s(m,r)x^r.
\]
If
\[
 F(\e^{-t})=\sum_{r\ge0}a_r\frac{t^r}{r!},
 \qquad q=1+h,
\]
then
\begin{equation}\label{eq:t-to-h-stirling}
 F(1+h)=\sum_{m\ge0}
 \left(\sum_{r=0}^{m}(-1)^r s(m,r)a_r\right)\frac{h^m}{m!}.
\end{equation}
Indeed,
\[
 \frac{(-\log(1+h))^r}{r!}
 =(-1)^r\sum_{m\ge r}s(m,r)\frac{h^m}{m!}.
\]
Thus a complete expansion in the geometrically natural variable $t$ can be
translated exactly to a literal Taylor expansion in $q-1$ by a finite
triangular transform at every order.

\section{Gaussian quick-reference coefficients}

Let
\[
 D=k(n-k),\qquad
 \Delta_{2r}=S_{2r}(n)-S_{2r}(k)-S_{2r}(n-k).
\]
The first three Faulhaber differences are
\begin{align*}
 \Delta_2={}&D(n+1),\\
 \Delta_4={}&D(n+1)\bigl(n(n+1)-D\bigr),\\
 \Delta_6={}&\frac{D(n+1)}2
 \bigl(2D^2-4Dn^2-5Dn+2n^4+4n^3+n^2-n\bigr).
\end{align*}
Consequently, for the centered normalized coefficient
\[
 H_{n,k}(t)=
 \e^{Dt/2}\frac{\qbinombase nk{\e^{-t}}}{\binom nk},
\]
we have
\begin{align*}
 \log H_{n,k}(t)
 ={}&\frac{\Delta_2}{24}t^2
     -\frac{\Delta_4}{2880}t^4
     +\frac{\Delta_6}{181440}t^6+\bigO(t^8),\\
 H_{n,k}(t)
 ={}&1+A_2t^2+\bigl(A_4+\tfrac12A_2^2\bigr)t^4\\
 &\quad+\bigl(A_6+A_2A_4+\tfrac16A_2^3\bigr)t^6+\bigO(t^8),
\end{align*}
where
\[
 A_2=\frac{\Delta_2}{24},\qquad
 A_4=-\frac{\Delta_4}{2880},\qquad
 A_6=\frac{\Delta_6}{181440}.
\]

In the literal coordinate $h=q-1$, the uncentered normalized polynomial has
\begin{align}
 \frac{\qbinombase nk{1+h}}{\binom nk}
 ={}&1+\frac D2h
 +\frac{D(3D+n-5)}{24}h^2
 +\frac{D(D-2)(D+n-3)}{48}h^3\notag\\
 &+\frac{D}{5760}\Bigl(
 15D^3+30D^2n-150D^2+5Dn^2-168Dn+487D\notag\\
 &\hspace{7.1em}{}-2n^3-4n^2+218n-500
 \Bigr)h^4+\bigO(h^5).
 \label{eq:gaussian-h-four}
\end{align}
Formula~\eqref{eq:gaussian-h-four} follows mechanically from
\eqref{eq:gaussian-log-all}, \eqref{eq:t-to-h-stirling}, and
\eqref{eq:log-product-recurrence}; it is included as a useful regression test
for symbolic implementations.

\chapter{Small-base and reciprocal coefficient tables}\label{app:small-base}

\section{A stable table for \texorpdfstring{$(aq;q)_\infty$}{the infinite aq product}}

The weighted distinct-partition expansion gives
\begin{align*}
 (aq;q)_\infty={}&1-aq-aq^2+a(a-1)q^3+a(a-1)q^4\\
 &+a(2a-1)q^5-a(a-1)^2q^6-a(a^2-3a+1)q^7\\
 &-a(a-1)(2a-1)q^8+\bigO(q^9).
\end{align*}
Multiplication by $1-a$ gives the corresponding expansion of
$(a;q)_\infty$.  The coefficient of $q^r$ is already stable in the finite
product $(aq;q)_{n-1}$ when $n>r$.

At $a=1$, Euler's pentagonal cancellation turns the table into
\[
 (q;q)_\infty
 =1-q-q^2+q^5+q^7+\bigO(q^{12}).
\]
At $a=-1$, it gives the generating product for partitions into distinct parts,
\[
 (-q;q)_\infty
 =1+q+q^2+2q^3+2q^4+3q^5+4q^6+5q^7+6q^8+\cdots.
\]

\section{Reciprocal transfer}

If a finite polynomial $P(q)=\sum_{r=0}^{D}c_rq^r$ is reciprocal,
$q^DP(q^{-1})=P(q)$, then the same coefficient list describes both ends:
\[
 P(q)=c_0+c_1q+\cdots+c_Dq^D
 =q^D\left(c_0+c_1q^{-1}+\cdots+c_Dq^{-D}\right).
\]
For Gaussian coefficients $c_r=p_{k,n-k}(r)$.  Hence every finite small-$q$
coefficient theorem immediately becomes a large-$q$ theorem, with the parity
factor $(-1)^D$ inserted on the negative real ray.

\chapter{Reproducible symbolic and numerical checks}\label{app:reproducibility}

\section{Companion program}

The archive contains \texttt{q\_expansion\_experiments.py}.  It requires only
Python, \texttt{sympy}, and \texttt{mpmath}; no specialized $q$-series package
is used.  Run it from the archive directory with
\begin{lstlisting}[style=pseudo]
python q_expansion_experiments.py
\end{lstlisting}
The program performs the following independent checks.

\begin{enumerate}[label=\textup{(\roman*)}]
\item It constructs $\qbinomq nk$ over $\Z[q]$ by the $q$-Pascal recurrence,
      checks the normalized expansion at $q=1$ through fourth order, and checks
      the value and derivative formula at $q=-1$, in each case for all
      $0\le k\le n\le20$.
\item It derives $\Delta_2$, $\Delta_4$, and $\Delta_6$ symbolically from
      Faulhaber sums and reduces them to polynomials in $n$ and
      $D=k(n-k)$.
\item It expands $\prod_{j\ge1}(1-aq^j)$ through $q^8$ by exact polynomial
      arithmetic.
\item It compares the Mellin expansion of $(q^x;q)_\infty$, the $q$-gamma
      expansion, the radial root-of-unity formula, and the double-scaling
      Gaussian formula with direct high-precision evaluation.
\end{enumerate}

\section{Recorded test instance}

With 80-decimal-digit working precision, the included run produced
\begin{center}\small
\begin{tabularx}{\textwidth}{@{}>{\raggedright\arraybackslash}p{0.19\textwidth}
  >{\raggedright\arraybackslash}p{0.43\textwidth}
  >{\raggedright\arraybackslash}X@{}}
\toprule
Test & Parameters and truncation & Absolute or signed error\\
\midrule
$(q^x;q)_\infty$ logarithm
 & $x=2.5$, $t=0.1$, through $t^6$
 & $2.64769017918096558\times10^{-14}$\\
$\log\Gamma_q(x)$
 & $x=3.7$, $t=0.08$, through $t^6$
 & $5.00835558895922029\times10^{-13}$\\
Primitive cubic root
 & $a=0.2$, $t=0.05$, through $t^6$
 & $3.4805992014908836\times10^{-11}$\\
Double scaling
 & $N=100$, $\alpha=0.4$, $\tau=1.7$, through $N^{-1}$
 & $5.34631320464741412\times10^{-8}$\\
Double scaling
 & same parameters, through $N^{-3}$
 & $-8.68801964159948774\times10^{-12}$\\
\bottomrule
\end{tabularx}
\end{center}
Each exact $q=1$ and $q=-1$ suite covered 230 index pairs and found no
discrepancy.  These checks do not replace the proofs, but they are sensitive to
signs, Bernoulli
normalizations, branch choices, and endpoint corrections that are easy to get
wrong in hand calculations.

\chapter{Compact formula map}\label{app:formula-map}

For reference, the principal expansion engines can be summarized in one page.

\begin{longtable}{@{}>{\raggedright\arraybackslash}p{0.20\textwidth}
                  >{\raggedright\arraybackslash}p{0.35\textwidth}
                  >{\raggedright\arraybackslash}p{0.37\textwidth}@{}}
\toprule
Regime & Natural normalization & Coefficient engine\\
\midrule
\endhead
Finite, regular $q=q_0$
 & Divide by the nonzero value at $q_0$
 & Logarithmic derivatives and Bell recurrence\\
Finite, resonant $q=q_0$
 & Extract $(q-q_0)^\nu$
 & Nonvanishing-factor logarithmic derivatives\\
$q\to1$, fixed degree
 & $q=\e^{-t}$; divide by the classical limit
 & Bernoulli numbers and Faulhaber differences\\
$q\to0$
 & None
 & Weighted restricted partitions\\
$q\to\infty$
 & Divide by the leading monomial
 & Reciprocity and the small-$q$ table\\
Finite root of unity
 & Extract the cyclotomic zero
 & Cyclotomic exponents, residue classes, Bell recurrence\\
Infinite $q\to1$, fixed $a$
 & Remove $\exp(-\Li_2(a)/t)$
 & Lambert series and polylogarithms\\
Infinite $q\to1$, $a=q^x$
 & Remove exponential, power, and gamma factors
 & Mellin residues and Bernoulli polynomials\\
Radial root of unity
 & Remove $\exp(-\Li_2(a^m)/(m^2t))$
 & Resonant residue class plus cyclic dilogarithm\\
$q=\e^{-\tau/N}$
 & Remove the order-$N$ dilogarithmic action
 & Singular Euler--Maclaurin; odd powers of $N^{-1}$\\
\bottomrule
\end{longtable}

\backmatter

\begin{thebibliography}{99}
\sloppy

\bibitem{ReshetnikovMonograph}
V.~Reshetnikov,
\emph{$q$-Pochhammer Symbols and $q$-Binomial Coefficients: Algebra,
Combinatorics, Analysis, Arithmetic, and Summation Theory},
ProveIt repository, August 2026,
\url{https://github.com/VladimirReshetnikov/ProveIt/blob/main/Analysis/FabiusFunction/docs/semi-formalized-research-frontiers/drafts/exponents-and-q-series/q_pochhammer_q_binomial_monograph/q_pochhammer_q_binomial_monograph.tex}.

\bibitem{McIntoshShifted}
R.~J. McIntosh,
``Some asymptotic formulae for $q$-shifted factorials,''
\emph{Ramanujan Journal} \textbf{3} (1999), 205--214.
\url{https://doi.org/10.1023/A:1006949508631}.

\bibitem{McIntoshHyper}
R.~J. McIntosh,
``Some asymptotic formulae for $q$-hypergeometric series,''
\emph{Journal of the London Mathematical Society} \textbf{51} (1995),
120--136.
\url{https://doi.org/10.1112/jlms/51.1.120}.

\bibitem{Moak}
D.~S. Moak,
``The $q$-analogue of Stirling's formula,''
\emph{Rocky Mountain Journal of Mathematics} \textbf{14} (1984), 403--413.
\url{https://doi.org/10.1216/RMJ-1984-14-2-403}.

\bibitem{GaroufalidisZagier}
S.~Garoufalidis and D.~Zagier,
``Asymptotics of Nahm sums at roots of unity,''
\emph{Ramanujan Journal} \textbf{55} (2021), 219--238.
\url{https://doi.org/10.1007/s11139-020-00266-x};
\url{https://arxiv.org/abs/1812.07690}.

\bibitem{ArdehaliRosengren}
A.~Arabi Ardehali and H.~Rosengren,
``A new product formula for $(z;q)_\infty$, with applications to asymptotics,''
arXiv:2602.11329, 2026.
\url{https://arxiv.org/abs/2602.11329}.

\bibitem{FantiniRella}
V.~Fantini and C.~Rella,
``Modular resurgence, $q$-Pochhammer symbols, and quantum operators from
mirror curves,''
\emph{Letters in Mathematical Physics} \textbf{116} (2026), article 39;
arXiv:2506.08265.
\url{https://arxiv.org/abs/2506.08265}.

\bibitem{ChengComanFantiniRella}
M.~C.~N. Cheng, I.~Coman, V.~Fantini, and C.~Rella,
``Modular resurgent structures for vectors,''
arXiv:2608.08902, 2026.
\url{https://arxiv.org/abs/2608.08902}.

\end{thebibliography}

\end{document}
